% !TeX program = pdflatex
% ======================================================================
% A Conditional Construction of Lorentzian Geometry from Euclidean Three-Space: Clifford Double-Cover Data, Gluing, and a Soldering Isomorphism
% Mathematical paper, dependency-ordered and theorem-centred.
%
% Source of mathematical truth for this draft:
% current Lean source in
% https://github.com/antaris82/The_Relational_Spin-Closure_Hypothesis
%
% Editorial attribution:
% Jan Seeck -- Editor
% ======================================================================

\documentclass[11pt,a4paper]{article}

\usepackage[margin=27mm]{geometry}
\usepackage[T1]{fontenc}
\usepackage[utf8]{inputenc}
\usepackage{microtype}
\usepackage{amsmath,amssymb,amsthm,mathtools}
\usepackage{enumitem}
\usepackage{booktabs}
\usepackage{array}
\usepackage{xcolor}
\usepackage{url}
\usepackage[
  hidelinks,
  pdftitle={A Conditional Construction of Lorentzian Geometry from Euclidean Three-Space: Clifford Double-Cover Data, Gluing, and a Soldering Isomorphism},
  pdfauthor={Jan Seeck (Editor)},
  pdfsubject={Conditional construction of Lorentzian geometry from Euclidean three-space using Clifford double-cover data, gluing data, and a soldering isomorphism}
]{hyperref}
\usepackage[nameinlink,noabbrev]{cleveref}

\setlength{\parindent}{1.4em}
\setlength{\parskip}{0pt}
\setlist[itemize]{topsep=.4em,itemsep=.2em}
\setlist[enumerate]{topsep=.4em,itemsep=.2em}

% ----------------------------------------------------------------------
% Theorem environments
% ----------------------------------------------------------------------

\newtheorem{theorem}{Theorem}[section]
\newtheorem{proposition}[theorem]{Proposition}
\newtheorem{lemma}[theorem]{Lemma}
\newtheorem{corollary}[theorem]{Corollary}

\theoremstyle{definition}
\newtheorem{definition}[theorem]{Definition}
\newtheorem{construction}[theorem]{Construction}
\newtheorem{hypothesis}[theorem]{Hypothesis}

\theoremstyle{remark}
\newtheorem{remark}[theorem]{Remark}

% ----------------------------------------------------------------------
% Notation
% ----------------------------------------------------------------------

\newcommand{\R}{\mathbb{R}}
\newcommand{\Ztwo}{\mathbb{Z}/2}
\newcommand{\Smodel}{\mathcal{S}}
\newcommand{\Cl}{\operatorname{Cl}}
\newcommand{\SpinG}{G_{\mathrm{Spin}}}
\newcommand{\LorG}{G_{\mathrm{Lor}}}
\newcommand{\End}{\operatorname{End}}
\newcommand{\GL}{\operatorname{GL}}
\newcommand{\id}{\operatorname{id}}
\newcommand{\im}{\operatorname{im}}
\newcommand{\kerop}{\operatorname{ker}}
\newcommand{\Ad}{\operatorname{Ad}}
\newcommand{\Int}{\operatorname{Int}}
\newcommand{\Diff}{\operatorname{Diff}}
\newcommand{\Cech}{\ifmmode\text{\v{C}ech}\else\v{C}ech\fi}
\newcommand{\cU}{\mathcal{U}}
\newcommand{\cR}{\mathcal{R}}
\newcommand{\cE}{\mathcal{E}}
\newcommand{\soB}{\mathfrak{so}(B)}
\newcommand{\wedgeS}{\Lambda^{2}\Smodel}
\newcommand{\dd}{\,\mathrm{d}}

\newcommand{\statusderived}{\textsc{derived}}
\newcommand{\statusdatum}{\textsc{additional datum}}
\newcommand{\statusopen}{\textsc{open}}

\DeclareMathOperator{\detop}{det}

% ----------------------------------------------------------------------
% Title
% ----------------------------------------------------------------------

\title{\textbf{A Conditional Construction of Lorentzian Geometry}\\[0.35em]
\large from Euclidean Three-Space: Clifford Double-Cover Data, Gluing, and a Soldering Isomorphism}

\author{Edited by Jan Seeck}
\date{13 September 2026}

\begin{document}
\maketitle

\begin{abstract}
We study a conditional construction whose local input is only the standard
positive-definite Euclidean space \(V=\R^{3}\).  From this datum one constructs
the real spin factor \(\Smodel=\R\oplus V\), its quadratic form
\(N(a,v)=a^{2}-\langle v,v\rangle\) of signature \((1,3)\), the real
Clifford algebra \(\Cl(V,q)\), a Clifford norm-one group \(\SpinG\), and a
surjective central two-fold covering
\[
   1\longrightarrow\{\pm1\}
   \longrightarrow \SpinG
   \overset{\rho}{\longrightarrow}\LorG
   \longrightarrow 1
\]
of the proper orthochronous Lorentz group \(\LorG\subset O(B)\) defined
below.  The \(B\)-skew-adjoint endomorphisms of \(\Smodel\) form the
six-dimensional orthogonal Lie algebra \(\mathfrak{so}(B)\), canonically
isomorphic to \(\Lambda^{2}\Smodel\).

These local constructions do not determine a global manifold.  We therefore
introduce explicit gluing data.  If the graph of the gluing equivalence
relation is closed, the index set is countable, and the transition maps are
smooth, the quotient is a Hausdorff second-countable smooth four-manifold
\(M\).

Independently, a \(\SpinG\)-valued \Cech\ 1-cocycle projects to a
Lorentz-valued cocycle.  This cocycle need not coincide with the tangent
transition cocycle \(D\phi_{ij}\).  Additional smooth soldering data
intertwine the two cocycles and determine a vector-bundle isomorphism from
the associated Lorentz vector bundle \(E_g\to M\) to \(TM\); the inverse is the
conventional \(E_g\)-valued solder form \(TM\to E_g\).  These data induce a smooth Lorentzian
metric on \(M\), orientation compatibility, and a continuous compatible
choice of the future timelike cone component in each tangent fibre.  They also imply that the
determinants of the tangent transition maps form a multiplicative \Cech\ 1-cocycle that is a continuous coboundary
and that the \(\{\pm1\}\)-valued \Cech\ 2-cocycle class associated with the
central extension is trivial on the chosen cover.  No identification of
these classes on the chosen cover with \(w_{1}(TM)\) or \(w_{2}(TM)\) is made.

Finally, on an explicit periodic gluing we study a one-parameter family of
\(\SpinG\)-valued \Cech\ 1-cocycles.  Soldering data exist for every finite
parameter.  Any two members are gauge equivalent under
\(\SpinG\)-valued \Cech\ 0-cochains; the corresponding sets of local soldering data
are explicitly bijective; and corresponding choices of soldering data induce the same
Lorentzian metric.  This result is restricted to that family and does not
construct a connection, holonomy, curvature, torsion, or dynamics.
\end{abstract}
\clearpage

\tableofcontents
\clearpage

% ======================================================================
\section{Introduction}
% ======================================================================

Standard Lorentzian spin geometry is usually formulated after
a smooth manifold, a Lorentz metric, and the associated frame bundle have
already been supplied; the Spin structure then appears as a lift of an
orthonormal frame bundle subject to a topological obstruction
\cite{LawsonMichelsohn,MilnorStasheff,Steenrod}.  Here the order of
dependence is deliberately reversed.  The starting object is only
\[
   V=\R^{3},\qquad
   \langle u,v\rangle=u_{1}v_{1}+u_{2}v_{2}+u_{3}v_{3},
\]
with no four-dimensional base space, tangent bundle, Lorentz form,
orientation, Spin bundle, connection, or curvature in the input.

There are then three logically separate construction questions.

First, what local Lorentzian and Clifford-algebraic structures can be
constructed from the local Euclidean algebra?  Second, what information must be
added before local copies of the resulting model can be assembled into a
global smooth manifold?  Third, after such a manifold has been obtained,
what additional datum is required to identify the Lorentz-valued cocycle with the tangent transition cocycle?

The separation of these questions is essential.  A local model is not a
manifold; fibre transition data do not determine which local points
are to be identified; and a Lorentz-valued cocycle is not, without
further structure, the derivative cocycle of a manifold atlas.  We make
each of these boundaries explicit and place the positive construction
theorems immediately after the relevant missing datum.

The local algebra belongs to the standard circle of ideas around real
Clifford algebras, Spin groups, spin factors and Lorentz cones
\cite{Chevalley,LawsonMichelsohn,JordanNeumannWigner,FarautKoranyi}.
The global steps use standard quotient, manifold, bundle and \Cech\ ideas
\cite{LeeSmooth,Husemoller,Steenrod,KobayashiNomizu}, but the point of the
present organization is not to assume the usual global Lorentzian structures
in advance.  Instead, the precise dependency chain is
\[
\boxed{
\begin{array}{c}
(\R^{3},\langle\cdot,\cdot\rangle)
\\[0.25em]\Downarrow\\[-0.15em]
(\Smodel,N,B,\Cl_{3},\SpinG,\rho,\soB)
\\[0.25em]
+\;\text{gluing data}
\\[-0.05em]\Downarrow\\[-0.15em]
\text{quotient covered by coordinate domains modelled on }\Smodel
\\[0.25em]
+\;\text{closed graph of the gluing relation + countable index + smooth transition maps}
\\[-0.05em]\Downarrow\\[-0.15em]
M^{4},\quad TM,\quad D\phi_{ij}
\\[0.25em]
+\;\text{soldering data}
\\[-0.05em]\Downarrow\\[-0.15em]
\begin{gathered}
(TM,g),\\
\text{orientation compatibility},\\
\text{compatible choice of a timelike cone component in each tangent fibre},\\
\text{tangent cocycle in the local frames defined by }A_i\text{ = projected Lorentz cocycle}
\end{gathered}
\end{array}}
\tag{1.1}
\]
The terms preceded by plus signs are additional data or hypotheses; they are not
derived from the preceding local algebra.

The following six statements organize the paper.

\begin{theorem}[Local construction]
\label{thm:intro-local}
From the specified Euclidean datum \(V=\R^{3}\), the construction produces:
\begin{enumerate}
\item a four-dimensional real spin factor \(\Smodel=\R\oplus V\) whose polarized quadratic form has signature \((1,3)\);
\item the real Clifford algebra \(\Cl_{3}=\Cl(V,q)\), with the paravector
      embedding of \(\Smodel\);
\item a Clifford norm-one group \(\SpinG\), the proper orthochronous Lorentz group \(\LorG\subset O(B)\) determined by the chosen future cone, and a continuous surjective
      homomorphism \(\rho:\SpinG\to\LorG\) whose central kernel is exactly
      \(\{\pm1\}\), together with continuous local sections;
\item the six-dimensional orthogonal Lie algebra \(\soB\) and a canonical linear
      isomorphism \(\Lambda^{2}\Smodel\simeq\soB\).
\end{enumerate}
\end{theorem}

\begin{theorem}[Globalization under explicit gluing hypotheses]
\label{thm:intro-global}
The local model does not determine a global quotient space.  Given explicit compatible gluing data \(B\) such that the graph of the gluing equivalence relation is closed, the index set is countable, and the transition maps are smooth, the quotient \(M_B\) is a Hausdorff second-countable smooth
four-manifold modelled on \(\Smodel\).
\end{theorem}

\begin{theorem}[Independence of tangent and Lorentz-valued cocycles]
\label{thm:intro-independence}
For the constructed smooth manifold, the tangent transition cocycle is \(D\phi_{ij}\).  Independently, a \(\SpinG\)-valued \Cech\ 1-cocycle projects to a Lorentz-valued cocycle \(\rho(\widetilde g_{ij})\).  Smooth gluing together with the latter cocycle does not force the two cocycles to agree.
\end{theorem}

\begin{theorem}[Soldering isomorphism]
\label{thm:intro-solder}
Soldering data are an additional smooth local family of fibrewise linear
isomorphisms satisfying the exact overlap intertwining law between
\(D\phi_{ij}\) and \(\rho(\widetilde g_{ij})\).
Conditional on their existence, they yield a vector-bundle isomorphism between \(E_g\), the Lorentz vector bundle defined by the cocycle, and
\(TM\), with continuous total-space maps in both directions and smooth local representatives.  They
also induce a smooth Lorentzian metric of signature \((1,3)\).
\end{theorem}

\begin{theorem}[Intertwining of tangent and Lorentz cocycles]
\label{thm:intro-compatibility}
For every choice of soldering data, the tangent transition cocycle expressed in the local frames defined by \(A_i\) equals the projected Lorentz-valued cocycle.  Consequently the determinants of the tangent transition maps form a multiplicative \Cech\ 1-cocycle that is a continuous coboundary, and the \(\{\pm1\}\)-valued \Cech\ 2-cocycle class for the central extension is trivial on the chosen cover.  The soldering data also yield a continuous compatible choice of the timelike cone component corresponding to \(C^{+}_{\mathrm{int}}\) in each tangent fibre.
\end{theorem}

\begin{theorem}[Gauge equivalence of the one-parameter family]
\label{thm:intro-deformation}
For the periodic gluing example and the one-parameter family of \(\SpinG\)-valued \Cech\ 1-cocycles considered in \cref{sec:deformation}, every finite parameter value
admits soldering data.  Any two parameter values define
cocycles gauge equivalent under \(\SpinG\)-valued \Cech\ 0-cochains.  The corresponding sets of local soldering data are explicitly bijective, and corresponding choices of soldering data induce the same
Lorentzian metric.
\end{theorem}

The hypotheses in \cref{thm:intro-global,thm:intro-solder} are not cosmetic.
They are exactly where the construction is no longer determined by the
initial Euclidean local data.  In particular, this paper does \emph{not}
prove an implication
\[
   \R^{3}\quad\Longrightarrow\quad\text{a uniquely determined smooth Lorentzian four-manifold}.
\]
It proves the local implication in \cref{thm:intro-local}, identifies the
additional global data required to continue, and then proves what follows
from those data.

% ======================================================================
\clearpage
\section{Euclidean input, the spin factor, and its Lorentzian quadratic form}
\label{sec:local}
% ======================================================================

\subsection{Euclidean input}

Throughout the local part of the paper,
\[
   V:=\R^{3}
\]
carries only its standard positive-definite bilinear form
\[
   \langle u,v\rangle
   :=u_{1}v_{1}+u_{2}v_{2}+u_{3}v_{3}.
\]
Set
\[
   q_{3}(v):=\langle v,v\rangle.
\]
No Lorentzian structure has yet been introduced.

\begin{definition}[Real spin factor]
\label{def:spin-factor}
Define
\[
   \Smodel:=\R\oplus V
\]
and, for \(x=(a,u)\), \(y=(b,v)\), define
\[
   x\circ y
   :=
   \bigl(ab+\langle u,v\rangle,\; av+bu\bigr).
\]
The distinguished unit is
\[
   \mathbf{1}_{\Smodel}:=(1,0).
\]
\end{definition}

The product is bilinear and commutative.  As an auxiliary local identity
for this chosen spin-factor construction one also has the Jordan identity
\[
   (x\circ y)\circ(x\circ x)
   =
   x\circ\bigl(y\circ(x\circ x)\bigr).
\]
Thus \((\Smodel,\circ,\mathbf{1}_{\Smodel})\) is the rank-two real spin
factor; compare the standard Jordan-algebra treatment in
\cite{JordanNeumannWigner,FarautKoranyi}.  This local identity is supporting
algebra, not an additional global construction claim.

\begin{definition}[Quadratic and bilinear forms]
\label{def:NSBS}
For \(x=(a,u)\in\Smodel\), set
\[
   N(x):=a^{2}-\langle u,u\rangle.
\]
Its polarization is
\[
   B(x,y)
   :=\frac{N(x+y)-N(x)-N(y)}{2}
   =ab-\langle u,v\rangle
\]
for \(y=(b,v)\).
\end{definition}

\begin{proposition}[Lorentz signature]
\label{prop:signature}
The form \(B\) is symmetric and nondegenerate, and has signature \((1,3)\).
In the standard basis
\[
   \mathbf{1}_{\Smodel},\quad
   (0,\varepsilon_{0}),\quad(0,\varepsilon_{1}),\quad(0,\varepsilon_{2}),
\]
its Gram matrix is
\[
   \operatorname{diag}(1,-1,-1,-1).
\]
\end{proposition}

\begin{proof}
The displayed formula
\(B((a,u),(b,v))=ab-\langle u,v\rangle\)
gives the Gram matrix directly.  Since the Euclidean form on \(V\) is
positive definite, the first basis direction is positive and the remaining
three are negative.  The determinant is nonzero, hence \(B\) is
nondegenerate.
\end{proof}

A second auxiliary local identity of the spin-factor construction is the
Cayley--Hamilton identity
\[
   x\circ x-\operatorname{tr}_{\Smodel}(x)\,x
   +N(x)\mathbf{1}_{\Smodel}=0,
   \qquad
   \operatorname{tr}_{\Smodel}(a,u)=2a.
\tag{2.1}
\]
The cone of squares
\[
   C_{\Smodel}:=
   \{x\in\Smodel:\exists y\in\Smodel,\ x=y\circ y\}
\tag{2.2}
\]
has the explicit coordinate description
\[
   C_{\Smodel}
   =\{(a,u)\in\Smodel:a\ge0,\ N(a,u)\ge0\}
   =\{(a,u)\in\Smodel:a\ge\lVert u\rVert\}.
\tag{2.2a}
\]
Thus \(C_{\Smodel}\) is the closed future Lorentz cone determined by the
scalar coordinate on \(\Smodel\).

\begin{definition}[Proper orthochronous Lorentz group]
\label{def:GLor}
Let \(\LorG\subset O(B)\) be the subgroup of linear automorphisms
\(F:\Smodel\to\Smodel\) satisfying
\[
   N(Fx)=N(x),\qquad
   x\in C_{\Smodel}\Longleftrightarrow Fx\in C_{\Smodel},
   \qquad
   \det F=1.
\tag{2.3}
\]
\end{definition}

The definition depends only on \((\Smodel,N,C_{\Smodel})\).  Relative to the
future cone \(C_{\Smodel}\), it is the proper orthochronous Lorentz group of
\(B\).  In particular,
no matrix realization over \(\mathbb C\), no pre-existing Lorentzian manifold,
and no frame bundle enters the construction.

\begin{remark}[Scope: local algebra only]
\label{bdry:no-manifold-yet}
At this point \(\Smodel\) is a four-dimensional real vector space with a
Lorentz form.  It is not a tangent space and there is no manifold,
chart, atlas, base point, or covering space in the construction.
\end{remark}

% ======================================================================
\clearpage
\section{Clifford algebra, a norm-one group, and a two-fold Lorentz covering}
\label{sec:clifford}
% ======================================================================

\subsection{The Euclidean Clifford algebra}

Let
\[
   \Cl_{3}:=\Cl(V,q_{3})
\]
be the real Clifford algebra characterized by the relation
\[
   \iota(v)^{2}=q_{3}(v)\mathbf{1}.
\tag{3.1}
\]
For the standard Euclidean basis
\(\varepsilon_{0},\varepsilon_{1},\varepsilon_{2}\) of \(V\), write
\(\gamma_{i}:=\iota(\varepsilon_{i})\).  Then
\[
   \gamma_{i}^{2}=1,\qquad
   \gamma_{i}\gamma_{j}=-\gamma_{j}\gamma_{i}
   \quad(i\neq j).
\tag{3.2}
\]
The signs are those of the \emph{positive} Euclidean quadratic form on
\(V\); the Lorentzian signature has already appeared in the distinct
four-dimensional spin factor \(\Smodel\).

\begin{definition}[Paravector embedding]
\label{def:paravector}
Define
\[
   A:\Smodel\longrightarrow\Cl_{3},
   \qquad
   A(a,v):=a\,\mathbf{1}+\iota(v).
\tag{3.3}
\]
\end{definition}

Let \(c\) denote the Clifford conjugation used in the construction.  Its
restriction to paravectors changes the sign of the vector part.  The key
identity is
\[
   A(x)c(A(x))=N(x)\mathbf{1}.
\tag{3.4}
\]

\begin{proposition}[Jordan--Clifford compatibility]
\label{prop:jordan-clifford}
For all \(x,y\in\Smodel\),
\[
   A(x\circ y)
   =
   \frac12\bigl(A(x)A(y)+A(y)A(x)\bigr).
\tag{3.5}\label{eq:jordan-clifford}
\]
Moreover, the image \(A(\Smodel)\) generates \(\Cl_{3}\) as a real
algebra.
\end{proposition}

\begin{proof}
Expand the product using
\[
   \iota(u)\iota(v)+\iota(v)\iota(u)
   =2\langle u,v\rangle\mathbf{1}.
\]
The scalar part becomes
\(2(ab+\langle u,v\rangle)\), while the vector part becomes
\(2(av+bu)\), proving \eqref{eq:jordan-clifford}.  Since \(A(0,v)=\iota(v)\), the image
contains every Clifford generator; the universal generation property of
the Clifford algebra then gives the second assertion.
\end{proof}

The Clifford algebra \(\Cl_{3}\) has its standard universal property:
any unital associative real algebra receiving a linear map \(f:V\to\mathcal A\) with
\(f(v)^{2}=q_{3}(v)\mathbf1\) receives a unique unital algebra homomorphism from
\(\Cl_{3}\); see \cite{Chevalley,LawsonMichelsohn}.

\subsection{The Clifford norm-one group}

\begin{definition}[Clifford norm-one group]
\label{def:spingroup}
An element \(g\in\Cl_{3}\) satisfies the Clifford norm-one condition when
\[
    g\,c(g)=c(g)\,g=1.
\tag{3.6}
\]
The corresponding units form a subgroup
\[
    \SpinG\subset\Cl_{3}^{\times}.
\tag{3.7}
\]
\end{definition}

For \(g\in\SpinG\), the Clifford action on paravectors preserves
\(A(\Smodel)\).  Transferring this action through the embedding \(A\)
defines a linear automorphism
\[
   \rho(g):\Smodel\longrightarrow\Smodel.
\]
For every \(g\in\SpinG\), the automorphism \(\rho(g)\) belongs to
\(\LorG\), and
\[
   \rho:\SpinG\longrightarrow\LorG
\tag{3.8}
\]
is a group homomorphism.  After the covering theorem below, this is the
paravector realization of the Lorentz double cover used throughout the paper.

\begin{theorem}[Two-fold central covering]
\label{thm:double-cover}
The homomorphism \(\rho\) is surjective and
\[
   \ker\rho=\{1,-1\}.
\tag{3.9}
\]
The kernel is central and has two elements.  With the topologies inherited
from the finite-dimensional Clifford model and the induced topology on
\(\LorG\), both groups are Hausdorff topological groups, \(\rho\) is
continuous, the kernel is discrete, and every \(k\in\LorG\) has an open
neighbourhood \(U\) admitting a continuous local section
\[
   s:U\longrightarrow\SpinG,
   \qquad
   \rho\circ s=\id_{U}.
\tag{3.10}
\]
\end{theorem}

\begin{proof}
We separate surjectivity, the kernel calculation, and the local-section
statement.

\emph{Surjectivity.}  Let \(F\in\LorG\) and put \(p=F(1_{\Smodel})\).
Because \(F\) preserves both \(N\) and the cone of squares,
\[
    N(p)=1,
    \qquad
    p\in C_{\Smodel},
    \qquad
    p_0\ge 0.
\]
Write \(p=(a,u)\).  From \(a^2-\langle u,u\rangle=1\) and \(a\ge0\) one
has \(a\ge1\).  Hence the following square root in the cone of squares
\[
 r:=\left(\sqrt{\frac{a+1}{2}},
      \frac{u}{2\sqrt{(a+1)/2}}\right)
\]
is well defined and satisfies \(r\circ r=p\) and \(N(r)=1\).  Set
\(\beta=A(r)\in\Cl_3\).  The Clifford norm identity makes \(\beta\) an element of \(\SpinG\), and the Jordan--Clifford compatibility gives
\[
   \rho(\beta)(1_{\Smodel})=p.
\]
Consequently
\[
   G:=\rho(\beta)^{-1}F
\]
fixes \(1_{\Smodel}\).  The stabilizer of \(1_{\Smodel}\) acts on the
spatial summand as an element \(R\in SO(3)\).  By Cartan--Dieudonn\'e,
\(R\) is a product of hyperplane reflections in unit spatial vectors.
For a unit vector \(u\), the Clifford generator \(\iota(u)\) acts on the
spatial summand as minus the corresponding hyperplane reflection.  Since
\(\det R=1\), an even number of reflections occurs; the accumulated minus
sign therefore disappears, and the product of the corresponding Clifford
generators is an element \(\gamma\in\SpinG\) with \(\rho(\gamma)=G\).  Thus
\(\rho(\beta\gamma)=F\).

\emph{Kernel.}  Suppose \(g\in\SpinG\) acts trivially on every paravector.
Evaluating at \(1_{\Smodel}\) first shows that Clifford reversion and
Clifford conjugation agree on \(g\).  The trivial-action equation then
implies
\[
    gA(x)=A(x)g
    \qquad\text{for every }x\in\Smodel.
\]
Because the paravector image generates \(\Cl_3\) as a real algebra, \(g\)
commutes with all of \(\Cl_3\).  The centre computed in the Clifford model is
\(\R\oplus\R\omega\), where
\(\omega=\gamma_0\gamma_1\gamma_2\) and \(\omega^2=-1\).  Hence
\(g=a+b\omega\).  For central elements Clifford conjugation fixes \(g\),
so the norm-one condition gives
\[
   (a+b\omega)^2=1
   \quad\Longrightarrow\quad
   a^2-b^2=1,
   \qquad
   2ab=0.
\]
The first equation excludes \(a=0\); therefore \(b=0\) and
\(a=\pm1\).  Conversely, \(\pm1\) act trivially.  Thus
\(\ker\rho=\{\pm1\}\); centrality and cardinality two follow immediately.

\emph{Continuous local sections.}  This construction uses only the Clifford model.  Define the linear map
\[
   \mathcal F(y)
   :=y+\sum_{i=0}^{2}\gamma_i y\gamma_i .
\]
In the eight-monomial Clifford frame this kills the six noncentral
monomials and multiplies the two central ones by \(4\), so
\(\mathcal F(y)\in\R\oplus\R\omega\).  Next define the auxiliary map
\[
   \mathcal R(F)
   :=A(F1_{\Smodel})
     +\sum_{i=0}^{2}A\!\left(F(0,\varepsilon_i)\right)\gamma_i .
\]
For \(g\in\SpinG\), direct Clifford multiplication gives
\[
   \mathcal R(\rho(g))
   =g\,\mathcal F(\operatorname{rev}g),
\]
so \(\mathcal R\) recovers \(g\) up to a central factor.  Identify the centre \(\R\oplus\R\omega\) with \(\mathbb C\) by the
real-algebra isomorphism
\[
   \jmath:\mathbb C\xrightarrow{\sim}\R\oplus\R\omega,
   \qquad \jmath(i)=\omega.
\]
Choose a real-linear map \(\pi_{\mathbb C}:\Cl_3\to\mathbb C\) whose
restriction to the centre is \(\jmath^{-1}\).  Define the continuous auxiliary scalar function
\[
   \Delta(F)
   :=\pi_{\mathbb C}\!\left(
      \mathcal R(F)\,c(\mathcal R(F))
   \right).
\]
The key identity is
\[
   \mathcal R(\rho(g))=\jmath(w)g
   \quad\Longrightarrow\quad
   \Delta(\rho(g))=w^2.
\]
Thus the auxiliary scalar function depends on the central factor through its square.  At the identity,
\(\Delta(1)=16\).  On the open neighbourhood where \(\Delta(F)\) lies in
a slit plane avoiding the branch cut, the principal complex square root is
continuous.  Dividing the Clifford element \(\mathcal R(F)\) by this square root
leaves exactly the two choices differing by \(\{\pm1\}\), and yields a
continuous \(\SpinG\)-valued section \(s\) with \(\rho(s(F))=F\) and \(s(1)=1\).
Finally, for arbitrary \(k\in\LorG\), choose one lift \(u_k\) of \(k\) and
left-translate the identity neighbourhood and its section.  This produces
an open neighbourhood of \(k\) with a continuous local section.  A
different choice of \(u_k\) changes that section only by the central
kernel element \(-1\).
\end{proof}

Thus one obtains the exact central extension
\[
   1\longrightarrow\{\pm1\}
   \longrightarrow\SpinG
   \overset{\rho}{\longrightarrow}\LorG
   \longrightarrow1.
\tag{3.11}\label{eq:spin-cover}
\]
The symbol \(\SpinG\) denotes the Clifford norm-one group of
Definition~\ref{def:spingroup}.  Its relation to \(\LorG\) is exactly the
central extension \eqref{eq:spin-cover}; no separate identification with the
even-subalgebra group \(\operatorname{Spin}(3)\subset\Cl_3^0\) is used.

A classical complex-matrix realization provides a comparison, but
it is not an input to \cref{thm:double-cover}.  In particular,
the construction does not define \(\SpinG\) by first declaring it to be
\(\mathrm{SL}(2,\mathbb C)\).

% ======================================================================
\clearpage
\section{The orthogonal Lie algebra and bivectors}
\label{sec:lie}
% ======================================================================

The local group construction has a parallel infinitesimal object that can
be defined without first making \(\SpinG\) into a Lie group.

\begin{definition}[Orthogonal Lie algebra]
\label{def:gB}
Let
\[
   \soB
   :=
   \bigl\{
      T\in\End_{\R}(\Smodel):
      B(Tx,y)+B(x,Ty)=0
      \ \text{for all }x,y\in\Smodel
   \bigr\}.
\tag{4.1}
\]
\end{definition}

For \(u,v\in\Smodel\), define
\[
   K_{u,v}(x)
   :=
   B(v,x)u-B(u,x)v.
\tag{4.2}\label{eq:Kuv}
\]
Then \(K_{u,v}\in\soB\), the map is alternating in \(u,v\), and it therefore
induces a linear map
\[
   \Phi:\Lambda^{2}\Smodel\longrightarrow\soB,
   \qquad
   u\wedge v\longmapsto K_{u,v}.
\tag{4.3}
\]

\begin{theorem}[Canonical identification \(\Lambda^2\Smodel\simeq\soB\)]
\label{thm:bivector}
The map \(\Phi\) is a linear isomorphism.  Consequently
\[
   \dim_{\R}\soB
   =
   \dim_{\R}\Lambda^{2}\Smodel
   =6.
\tag{4.4}
\]
Moreover, \(\soB\) is closed under the commutator
\[
   [A,C]:=AC-CA,
\tag{4.5}
\]
so it is a six-dimensional real Lie algebra.
\end{theorem}

\begin{proof}
Skewness of \(K_{u,v}\) follows by expanding \eqref{eq:Kuv} and using the
symmetry of \(B\).  Closure of \(\soB\) under the commutator follows by
applying the skewness identity twice:
\[
  B(ACx,y)=-B(Cx,Ay)=B(x,CAy),
\]
and similarly with \(A\) and \(C\) interchanged.

It remains to prove that \(\Phi\) is bijective.  Let
\[
   \tau=(1,0),
   \qquad
   b_i=(0,\varepsilon_i),\quad i=0,1,2,
\]
where \(\varepsilon_0,\varepsilon_1,\varepsilon_2\) is the fixed Euclidean
basis of \(V\).  For \(A\in\soB\), introduce the six coefficients
\[
\begin{aligned}
   c_0&=(A\tau)_V^{0},&
   c_1&=(A\tau)_V^{1},&
   c_2&=(A\tau)_V^{2},\\
   c_3&=\langle (Ab_0)_V,\varepsilon_1\rangle,&
   c_4&=\langle (Ab_0)_V,\varepsilon_2\rangle,&
   c_5&=\langle (Ab_1)_V,\varepsilon_2\rangle.
\end{aligned}
\tag{4.6a}\label{eq:gb-six-coordinates}
\]
The \(B\)-skew condition determines every other matrix coefficient from
these six: it forces \((A\tau)_0=0\), determines the scalar component of
\(A(0,u)\) by \(\langle(A\tau)_V,u\rangle\), and makes the spatial block
antisymmetric.  Hence if all six numbers in
\eqref{eq:gb-six-coordinates} vanish, then \(A\tau=0\) and
\(Ab_i=0\) for \(i=0,1,2\).  Since \(\tau,b_0,b_1,b_2\) span
\(\Smodel\), one has \(A=0\).  Thus the six coefficients determine \(A\) uniquely.

Conversely, for \(c=(c_0,\ldots,c_5)\in\R^6\), define the explicit
bivector
\[
\begin{aligned}
   z(c):={}&-c_0\,\tau\wedge b_0
            -c_1\,\tau\wedge b_1
            -c_2\,\tau\wedge b_2\\
          &+c_3\,b_0\wedge b_1
            +c_4\,b_0\wedge b_2
            +c_5\,b_1\wedge b_2 .
\end{aligned}
\tag{4.6b}\label{eq:zmap}
\]
A direct evaluation of \(K_{u,v}\) on the four basis vectors shows that
the six coordinates of \(\Phi(z(c))\) are exactly \(c_0,\ldots,c_5\).
For any \(A\in\soB\), taking \(c\) to be the six coordinates of \(A\)
therefore gives an endomorphism \(\Phi(z(c))\) with the same six
coordinates as \(A\); uniqueness of the six coefficients yields \(\Phi(z(c))=A\).  Thus \(\Phi\) is
surjective.

For injectivity, the map \(z:\R^6\to\Lambda^2\Smodel\) in
\eqref{eq:zmap} is injective because the six coordinates of
\(\Phi(z(c))\) recover \(c\).  Since
\[
   \dim\R^6=6=\dim\Lambda^2\Smodel,
\]
this injective linear map is also surjective.  If \(\Phi(z)=\Phi(z')\),
write \(z=z(c)\) and \(z'=z(c')\); applying the six coordinate functionals
gives \(c=c'\), hence \(z=z'\).  Therefore \(\Phi\) is injective as well,
and so it is a linear isomorphism.  The dimension statement follows.
\end{proof}

\begin{remark}[Open problem: relation with the Lie algebra of \(\SpinG\)]
\label{bdry:lie-spin}
The statement
\[
   \operatorname{Lie}(\SpinG)\simeq\soB
\tag{4.7}
\]
is not part of the present construction.  The group \(\SpinG\) is
established here as a topological group, not as a smooth Lie group.
Accordingly, the differential
\[
   d\rho_{e}:\operatorname{Lie}(\SpinG)
   \longrightarrow\soB
\tag{4.8}
\]
has not been constructed.  These statements remain open in the present paper.
\end{remark}

% ======================================================================
\clearpage
\section{From local model to global quotient space}
\label{sec:globalization}
% ======================================================================

The next step is deliberately not phrased as a consequence of the Spin
group.  The reason is structural: a fibre transition tells us how to
compare fibres over a point already known to lie in two patches.  It does
not say which points of two local copies represent the same point of a
global quotient space.

\subsection{Gluing data}

Let \(I\) be an index set.  A local piece is an open subset
\(D_i\subset\Smodel\).

\begin{definition}[Gluing datum]
\label{def:basegluing}
A gluing datum \(B\) consists of:
\begin{enumerate}
\item open sets \(D_i\subset\Smodel\);
\item open overlap domains \(W_{ij}\subset D_i\);
\item continuous maps
      \(\phi_{ij}:\Smodel\to\Smodel\), used on \(W_{ij}\);
\end{enumerate}
such that, on the indicated domains,
\[
\begin{aligned}
   &W_{ii}=D_i,
   &&\phi_{ii}(x)=x,\\
   &\phi_{ij}(W_{ij})\subset W_{ji},
   &&\phi_{ji}(\phi_{ij}(x))=x,
\end{aligned}
\tag{5.1}
\]
and whenever \(x\in W_{ij}\) and
\(\phi_{ij}(x)\in W_{jk}\),
\[
   x\in W_{ik},
   \qquad
   \phi_{jk}(\phi_{ij}(x))=\phi_{ik}(x).
\tag{5.2}
\]
\end{definition}

This is an explicit \emph{sufficient} gluing datum.  We do not
claim that every field in Definition~\ref{def:basegluing} is separately necessary,
or that this is the unique minimal encoding of overlap information.

Form the tagged disjoint union
\[
   X_{0}:=\coprod_{i\in I}D_i
\tag{5.3}
\]
and define
\[
   (i,x)\sim(j,y)
   \quad\Longleftrightarrow\quad
   x\in W_{ij}\ \text{ and }\ \phi_{ij}(x)=y.
\tag{5.4}
\]
The axioms above make \(\sim\) an equivalence relation.  Define
\[
   M_B:=X_{0}/\!\sim
\tag{5.5}
\]
with the quotient topology.

\begin{proposition}[Topological quotient construction]
\label{prop:top-reconstruct}
The quotient map \(X_{0}\to M_B\) is open.  Each local piece \(D_i\)
embeds homeomorphically onto an open subset of \(M_B\), and these images
cover \(M_B\).  Thus the images of the \(D_i\) form a topological atlas of open coordinate domains modelled on \(\Smodel\), before any smoothness assertion is made.
\end{proposition}

\begin{proof}
The saturation of an open subset of \(X_{0}\) is open because each
overlap set \(W_{ij}\) is open and \(\phi_{ij}\) is a homeomorphism
between \(W_{ij}\) and \(W_{ji}\), with inverse \(\phi_{ji}\).  Hence the
quotient map is open.  Equality of two points from the same tagged piece
forces equality of their local coordinates by \(\phi_{ii}=\id\), so the
map from each \(D_i\) is injective.  Openness and continuity then give an
open embedding.  Surjectivity of the quotient map proves that the chart
images cover.
\end{proof}

Hausdorffness is not a formal consequence of Definition~\ref{def:basegluing}; the
usual line-with-two-origins phenomenon already shows why a separation
condition is needed.

\begin{definition}[Closed gluing relation]
\label{def:closedgraph}
The equivalence relation associated with \(B\) has closed graph when
\[
   \Gamma_B
   :=
   \{(p,q)\in X_0\times X_0:p\sim q\}
\tag{5.6}
\]
is closed.
\end{definition}

\begin{proposition}[Hausdorffness and second countability]
\label{prop:haus-second}
If \(\Gamma_B\) is closed, then \(M_B\) is Hausdorff.  If, in addition,
\(I\) is countable, then \(M_B\) is second countable.
\end{proposition}

\begin{proof}
Let \(q:X_0\to M_B\) be the quotient map.  By
Proposition~\ref{prop:top-reconstruct}, \(q\) is continuous, open, and surjective.
Consequently
\[
   q\times q:X_0\times X_0\longrightarrow M_B\times M_B
\]
is again an open quotient map.  The preimage of the diagonal is exactly the
gluing relation:
\[
   (q\times q)^{-1}(\Delta_{M_B})
   =\{(p,q):p\sim q\}
   =\Gamma_B.
\]
If \(\Gamma_B\) is closed, the quotient property therefore implies that
\(\Delta_{M_B}\) is closed.  This is equivalent to Hausdorffness.

For second countability, each \(D_i\) is an open subspace of the
finite-dimensional real vector space \(\Smodel\), hence is second
countable.  If \(I\) is countable, the tagged disjoint union
\(X_0=\coprod_i D_i\) is second countable.  An open quotient of a
second-countable space is second countable: the images under \(q\) of a
countable basis form a countable basis of \(M_B\).  Hence \(M_B\) is second
countable.
\end{proof}

\subsection{Why the quotient space requires additional data}

The need for gluing data is itself a theorem, not merely a warning.

\begin{theorem}[Local pieces do not determine the quotient space]
\label{thm:base-nonderivability}
There exist two gluing data over the same two indexed local domains
such that the local pieces agree exactly but their quotient spaces are not
homeomorphic.
\end{theorem}

\begin{proof}
Take a nonempty connected open local domain \(D\subset\Smodel\) and two
copies indexed by \(\{\mathsf{false},\mathsf{true}\}\).  In the first
gluing, set \(W_{ij}=D\) for every pair of indices and use the identity
identification.  In the second, take \(W_{ij}=\varnothing\) for \(i\neq j\)
and again use the identity where defined.  The first quotient is connected;
the second is the disjoint union of two nonempty open components and is not
connected.  Hence no homeomorphism can exist between them.
\end{proof}

\begin{corollary}
\label{cor:fibre-not-base}
No datum that only prescribes a fibre transition law over already specified
overlaps can, by itself, determine the missing identifications between distinct
local pieces.
\end{corollary}

\subsection{Smooth gluing}

For the constructed atlas, the coordinate change from piece \(i\) to
piece \(j\) is precisely \(\phi_{ij}\) on \(W_{ij}\).

\begin{definition}[Smooth gluing]
\label{def:smoothgluing}
A gluing datum is smoothly compatible when
\[
   \phi_{ij}|_{W_{ij}}
\]
is \(C^\infty\) for every \(i,j\).
\end{definition}

\begin{theorem}[Conditional smooth four-manifold]
\label{thm:smooth-manifold}
Suppose:
\begin{enumerate}
\item \(B\) is a gluing datum over \(\Smodel\);
\item its gluing relation is closed;
\item the index set \(I\) is countable;
\item \(B\) satisfies Definition~\ref{def:smoothgluing}.
\end{enumerate}
Then \(M_B\) is Hausdorff, second countable, and a \(C^\infty\)
four-manifold modelled on \(\Smodel\).
\end{theorem}

\begin{proof}
The Hausdorff and second-countable assertions are exactly
Proposition~\ref{prop:haus-second}.  It remains to identify the smooth structure of
the actual quotient atlas.

For each nonempty piece \(D_i\), the open embedding of
Proposition~\ref{prop:top-reconstruct} yields a chart whose target is \(D_i\).  On an
overlap of the \(i\)- and \(j\)-charts, the source of the coordinate change
is not merely contained in the specified overlap set: it is exactly
\(W_{ij}\).  If \(y\in W_{ij}\), the quotient relation identifies
\((i,y)\) with \((j,\phi_{ij}(y))\), and therefore the coordinate change is
literally
\[
     y\longmapsto\phi_{ij}(y).
\]
Thus every transition map of the constructed atlas is a restriction of
one of the functions already required to be \(C^\infty\) by
Definition~\ref{def:smoothgluing}.  The transition maps are therefore smooth, so the resulting atlas is a smooth atlas and \(M_B\) is a smooth manifold.
Finally, the model space is
\(\Smodel=\R\oplus\R^3\), hence \(\dim_{\R}\Smodel=4\).

Notice that the argument proves sufficiency of the smooth-gluing
hypothesis for this constructed atlas.  It does not claim the converse or
necessity of every specified smoothness condition.
\end{proof}

\begin{proposition}[Counterexample: topological gluing need not be smooth]
\label{prop:nonsmooth-gluing}
There are gluing data satisfying the topological gluing axioms whose
transition is a homeomorphism but not differentiable at the origin.  An explicit example is
\[
   H(t,x)_0=t,
   \qquad
   H(t,x)_k=x_k+|t| \quad(k=1,2,3),
\]
with inverse
\[
   H^{-1}(t,x)_0=t,
   \qquad
   H^{-1}(t,x)_k=x_k-|t|.
\]
Both maps are continuous, so the topological quotient construction applies.  But
along the line \(t\mapsto(t,0)\), every shifted spatial component contains
\(|t|\), which is not differentiable at \(t=0\).  Hence the smooth-gluing
condition fails.
\end{proposition}

\begin{remark}[Sufficiency, not necessity]
\label{rem:smooth-sufficient}
For the quotient construction used here, smooth gluing is proved
sufficient.  A converse is not asserted.  In particular, quotient
representative selection may cause the constructed atlas to omit some
index-labelled chart even when the corresponding specified
\(\phi_{ij}\) is non-smooth.
\end{remark}

% ======================================================================
\clearpage
\section{\Cech\ 1-cocycles for the Clifford norm-one group and tangent transitions}
\label{sec:tangent}
% ======================================================================

Fix from now on a smooth gluing \(B\) satisfying the hypotheses needed
to obtain \(M=M_B\).

\subsection{Tangent transitions}

Because the actual atlas transition is \(\phi_{ij}\), the tangent bundle
transition is its derivative.

\begin{proposition}[Tangent transition law]
\label{prop:tangent-transition}
For \(y\in W_{ij}\), the tangent coordinate change of the
constructed manifold is
\[
   T_{ij}(y):=D\phi_{ij}(y):
   \Smodel\longrightarrow\Smodel.
\tag{6.1}\label{eq:tangent-transition}
\]
It is a linear isomorphism, with inverse
\[
   D\phi_{ji}(\phi_{ij}(y)).
\tag{6.2}\label{eq:tangent-inverse}
\]
The transitions satisfy the derivative cocycle law by the chain rule.
\end{proposition}

The invertibility in \eqref{eq:tangent-inverse} is derived from
\(\phi_{ji}\circ\phi_{ij}=\id\) on the open overlap domain rather than
inserted as a separate hypothesis.

\begin{remark}[Local coordinate identifications do not identify the cocycles]
Although each chart identifies tangent vectors locally with the model
space \(\Smodel\), these local identifications do not identify the tangent
transition cocycle with the independent Lorentz cocycle.  The relevant
transition law is \eqref{eq:tangent-transition}.
\end{remark}

\subsection{A group-valued \Cech\ 1-cocycle}

Let \(\cU=\{U_i\}_{i\in I}\) be the cover of \(M\) by the chart images.
A \(\SpinG\)-valued \Cech\ 1-cocycle consists of continuous maps
\[
   \widetilde g_{ij}:U_i\cap U_j\longrightarrow\SpinG
\tag{6.3}
\]
satisfying the \Cech\ cocycle law
\[
   \widetilde g_{ij}(x)\widetilde g_{jk}(x)
   =
   \widetilde g_{ik}(x)
\tag{6.4}
\]
on triple overlaps.  Projection gives the Lorentz-valued cocycle
\[
   g_{ij}:=\rho\circ\widetilde g_{ij}:
   U_i\cap U_j\longrightarrow\LorG.
\tag{6.5}
\]
This is a fibre transition cocycle on the same cover, but it is not the
tangent transition cocycle by definition.

\begin{theorem}[The tangent and Lorentz transition cocycles need not coincide]
\label{thm:transition-independence}
There exist a smooth two-chart gluing \(B\) and a \(\SpinG\)-valued \Cech\ 1-cocycle \(\widetilde g\) such that
\[
   g_{01}(x)=\id_{\Smodel}
   \quad\text{for every }x,
\tag{6.6}
\]
while
\[
   T_{01}(y)=2\,\id_{\Smodel}
   \quad\text{for every }y.
\tag{6.7}
\]
Hence the gluing data and the \(\SpinG\)-valued cocycle do not force
\(T_{ij}=g_{ij}\).
\end{theorem}

\begin{proof}
Take two copies of the whole local model with
\[
   \phi_{01}(y)=2y,\qquad
   \phi_{10}(y)=\tfrac12 y.
\]
All overlap domains are the whole model, and the gluing is smooth.
Thus \(D\phi_{01}=2\,\id\).  Independently, choose the pointwise trivial \(\SpinG\)-valued \Cech\ 1-cocycle \(\widetilde g_{ij}=1\); its projection is the identity.
Evaluating at \(\mathbf1_{\Smodel}\) distinguishes the two transformations.
\end{proof}

\begin{remark}
The conclusion is one of non-derivability, not non-existence.  The same rescaling example admits compatible soldering data.  Equality of transition matrices is therefore not the correct invariant criterion for
coupling the two bundles.
\end{remark}

% ======================================================================
\clearpage
\section{A soldering isomorphism between the Lorentz vector bundle and the tangent bundle}
\label{sec:solder}
% ======================================================================

Let \(E_g\to M\) denote the rank-four real vector bundle defined by the
Lorentz-valued \Cech\ 1-cocycle \(g_{ij}\).  The additional object considered in
this section is a fibrewise linear identification \(s:E_g\to TM\) whose local
representatives vary smoothly and intertwine the two transition cocycles.  The
inverse map \(\theta:=s^{-1}:TM\to E_g\) has the conventional direction of an
\(E_g\)-valued solder form \cite{KobayashiNomizu,Nakahara}.

\subsection{Pointwise fibre isomorphisms}

At the weakest level one may choose, for each chart and each point of its image in \(M\), a
linear isomorphism
\[
   e_i(x):\Smodel\xrightarrow{\;\sim\;}T_xM
\tag{7.1}
\]
such that on \(U_i\cap U_j\),
\[
   e_j(x)v
   =
   e_i(x)\bigl(g_{ij}(x)v\bigr).
\tag{7.2}\label{eq:weak-solder-law}
\]
This algebraic condition is sufficient to transport the Lorentz form \(B\) to each tangent fibre, but it imposes no regularity on the field
\(x\mapsto e_i(x)\).

\begin{proposition}[Counterexample: pointwise fibre isomorphisms need not be continuous]
\label{prop:weak-solder}
There exist pointwise families of fibre isomorphisms obeying \eqref{eq:weak-solder-law} whose chart
representative is discontinuous.  Consequently pointwise fibre isomorphisms alone do not produce a continuous vector-bundle isomorphism or a smooth Lorentzian metric.
\end{proposition}

\begin{proof}[Explicit construction]
Take the one-chart quotient with identity gluing and the trivial \(\SpinG\)-valued \Cech\ 1-cocycle.  In this case \eqref{eq:weak-solder-law} imposes no
nontrivial overlap constraint.  Let \(x_0\) be the chart point represented
by the zero vector and define the frame field
\[
 e(x)=
 \begin{cases}
   2\,\id_{\Smodel}, & x=x_0,\\
   \id_{\Smodel}, & x\ne x_0.
 \end{cases}
\tag{7.2a}\label{eq:bad-solder}
\]
Every fibre map is a linear isomorphism, so this is a valid pointwise family of fibre isomorphisms.
Its local representative, however, jumps from \(\id\) to \(2\id\) at
\(x_0\), and is therefore discontinuous.  The discontinuity is also visible in
the transported quadratic form: evaluated on the distinguished unit, the
coefficient is \(1\) away from \(x_0\), whereas at \(x_0\) the inverse
frame contributes a factor \(1/2\), giving \(1/4\).  Thus even the induced
fibrewise metric coefficient is discontinuous.  This discontinuous pointwise frame field is therefore not induced by soldering data of Definition~\ref{def:soldering}.
\end{proof}

\subsection{Local representatives of a soldering isomorphism}

We describe the soldering isomorphism by local representatives in the chart coordinates.  Let
\[
   A_i(y)\in\GL(\Smodel)
\]
be a continuous linear automorphism for \(y\in D_i\).

\begin{definition}[Local soldering data]
\label{def:soldering}
Local soldering data for \((B,\widetilde g)\) are a family
\[
   A_i:D_i\longrightarrow\GL(\Smodel)
\tag{7.3}
\]
such that:
\begin{enumerate}
\item the map \(y\mapsto A_i(y)\), viewed in the finite-dimensional space
      of continuous linear maps, is \(C^\infty\) on \(D_i\);
\item on \(y\in W_{ij}\),
\[
  A_j(\phi_{ij}(y))
  =
  T_{ij}(y)\circ A_i(y)\circ
  g_{ij}(x_y),
\tag{7.4}\label{eq:soldering-law}
\]
where \(x_y\in M\) is the quotient point represented by \((i,y)\).
\end{enumerate}
\end{definition}

Equation \eqref{eq:soldering-law} is the local intertwining identity.
The tangent map \(T_{ij}\) converts \(i\)-chart tangent components to
\(j\)-chart tangent components, while the \Cech\ convention for \(E_g\) identifies components in the two local
trivializations by \(g_{ij}\).  The equation is therefore the compatibility
condition between the two transition cocycles.

The inverse family \(A_i(y)^{-1}\) is automatically smooth; no separate smoothness assumption for the inverse is needed.

\begin{theorem}[Soldering isomorphism]
\label{thm:bundle-equivalence}
The local soldering data determine a vector-bundle isomorphism
\[
   s:E_g\longrightarrow TM
\]
covering \(\id_M\), at the following regularity level:
\begin{enumerate}
\item continuous total-space map and continuous inverse;
\item linear isomorphism on each fibre;
\item \(C^\infty\) local representatives \(A_i\);
\item \(C^\infty\) local representatives \(A_i^{-1}\) for the inverse;
\item exact intertwining of the Lorentz-valued and tangent transition cocycles.
\end{enumerate}
Conversely, a vector-bundle isomorphism with these properties yields local soldering data of Definition~\ref{def:soldering}.
\end{theorem}

\begin{proof}
Write \(E_g\) in the local trivializations determined by
the cocycle \(g_{ij}\), and write \(TM\) in the tangent trivializations of
the constructed atlas.  On the \(i\)-th chart define the local total-space
map
\[
   (y,v)\longmapsto (y,A_i(y)v).
\tag{7.4a}\label{eq:local-bundle-map}
\]
If \(y\in W_{ij}\), changing the source trivialization by \(g_{ij}\) and
the target trivialization by \(T_{ij}\) gives, by
\eqref{eq:soldering-law}, the same point as first applying
\eqref{eq:local-bundle-map} in chart \(i\).  Hence the local maps agree on
equivalent representatives and descend to a globally defined map of total
spaces covering \(\id_M\).  Each fibre map is a linear isomorphism because
each \(A_i(y)\) is.

The local representative of the descended map in the \(i\)-th source and
target trivializations is exactly \(A_i\).  Since \(A_i\) is smooth, it is
continuous, and the standard local-trivialization criterion gives
continuity of the total-space map.  Pointwise inversion on
\(\GL(\Smodel)\) is smooth; hence \(y\mapsto A_i(y)^{-1}\) is smooth.  The
same argument applied to these inverse representatives yields a continuous
inverse total-space map.  This proves the five stated properties.

Conversely, suppose such a vector-bundle isomorphism is given.  Read
its local representative in the chart trivializations and call it
\(A_i(y)\).  Fibrewise linear isomorphism gives
\(A_i(y)\in\GL(\Smodel)\), smoothness of the local representatives gives \(C^\infty\)-regularity,
and equality of the globally defined map when read in the two overlapping
trivializations is precisely the intertwining identity
\eqref{eq:soldering-law}.  Thus the local representatives form local soldering data of the stated kind.  The two constructions are inverse at the level asserted
in the theorem.
\end{proof}

\begin{remark}[Exact regularity level]
We establish a vector-bundle isomorphism with continuous total-space maps in both
directions, fibrewise linear isomorphisms, and smooth local representatives
for the map and its inverse.  No smooth structure on the total spaces is used to state a stronger smooth bundle isomorphism.  In particular, the limitation concerns the level at which the
global equivalence is stated, not the smoothness of the projected Lorentz
transition representatives, which is forced in chart coordinates by the soldering data.
\end{remark}

Indeed, \eqref{eq:soldering-law} can be solved for \(g_{ij}\):
\[
   g_{ij}(x_y)
   =
   A_i(y)^{-1}\circ T_{ij}(y)^{-1}
   \circ A_j(\phi_{ij}(y)),
\tag{7.5}
\]
up to the convention used in \eqref{eq:soldering-law}.  Hence the soldering data force
the projected Lorentz transition to possess smooth local
representatives in chart coordinates even though the underlying \(\SpinG\)-valued \Cech\ 1-cocycle is only assumed continuous.

\subsection{Gauge freedom of the soldering isomorphism}

For fixed \(B\) and fixed \(\SpinG\)-valued \Cech\ 1-cocycle, the soldering data need not be unique.  Let \(\operatorname{Gau}(TM)\) denote the group of smooth vector-bundle automorphisms of \(TM\) covering \(\id_M\).  It acts on the set of local soldering data.

\begin{theorem}[Torsor of soldering isomorphisms]
\label{thm:solder-torsor}
Whenever the set of soldering isomorphisms represented by local data for fixed \((B,\widetilde g)\)
is nonempty, the action of \(\operatorname{Gau}(TM)\) on it is free and transitive.
\end{theorem}

\begin{proof}
Let \(A=\{A_i\}\) and \(A'=\{A_i'\}\) be two choices of local soldering data for the same
gluing datum and the same \(\SpinG\)-valued \Cech\ 1-cocycle.  Define
\[
   a_i(y):=A_i'(y)\,A_i(y)^{-1}.
\tag{7.6}\label{eq:gauge-between-solders}
\]
The field \(a_i\) is smooth because multiplication and inversion in the
finite-dimensional general linear group are smooth.  Applying the intertwining identity \eqref{eq:soldering-law} once to \(A\) and once to \(A'\), and cancelling the common Lorentz-valued transition, gives
\[
   a_j(\phi_{ij}(y))\,T_{ij}(y)
   =T_{ij}(y)\,a_i(y),
\tag{7.7}\label{eq:tangent-gauge-law}
\]
which is exactly the gauge-compatibility law for \(TM\).  By
construction,
\[
   (a\cdot A)_i(y)=a_i(y)A_i(y)=A_i'(y),
\]
so the action is transitive.

For freeness, if \(a\cdot A=a'\cdot A\), then
\(a_i(y)A_i(y)=a_i'(y)A_i(y)\) at every point.  Right cancellation by the
invertible map \(A_i(y)\) gives \(a_i(y)=a_i'(y)\) for all \(i,y\), hence
\(a=a'\).  Pointwise composition and inversion preserve
\eqref{eq:tangent-gauge-law}, so these fields form a group and the action
is free and transitive.
\end{proof}

Thus existence is an additional requirement, while uniqueness is neither
expected nor claimed.

% ======================================================================
\clearpage
\section{The induced Lorentzian metric}
\label{sec:metric}
% ======================================================================

Let \(A=\{A_i\}\) be a choice of soldering data.

\subsection{Metric coefficients}

In chart \(i\), define
\[
   g_i(y)(u,v)
   :=
   B\bigl(A_i(y)^{-1}u,\;A_i(y)^{-1}v\bigr).
\tag{8.1}\label{eq:metric-coeff}
\]

\begin{theorem}[Induced Lorentzian metric]
\label{thm:smooth-metric}
The fields \(g_i\) satisfy:
\begin{enumerate}
\item for fixed \(u,v\in\Smodel\), the function
      \(y\mapsto g_i(y)(u,v)\) is \(C^\infty\) on \(D_i\);
\item \(g_i(y)\) is symmetric and nondegenerate;
\item on every overlap domain,
\[
  g_j(\phi_{ij}(y))
     \bigl(T_{ij}(y)u,T_{ij}(y)v\bigr)
  =
  g_i(y)(u,v);
\tag{8.2}\label{eq:metric-transform}
\]
\item in the \(A_i(y)\)-frame, the form has signature \((1,3)\).
\end{enumerate}
Hence the \(g_i\) are the local coefficient fields of a smooth section
\[
   g\in\Gamma(\operatorname{Sym}^{2}T^{*}M)
\]
that is nondegenerate of signature \((1,3)\) at every point.  Thus \(g\) is a smooth
Lorentzian metric on \(M\) in the standard semi-Riemannian sense
\cite{ONeill,Nakahara}.
\end{theorem}

\begin{proof}
Smoothness follows from smoothness of \(A_i^{-1}\) and bilinearity of
\(B\).  Symmetry and nondegeneracy are inherited from \(B\) through the
linear isomorphism \(A_i(y)\).  For the transformation law, use
\eqref{eq:soldering-law} to express the inverse \(j\)-frame after application of
\(T_{ij}\) in terms of the inverse \(i\)-frame and the projected Lorentz-valued transition.  Since \(g_{ij}(x)\in\LorG\) preserves \(B\), these
Lorentz factors cancel and yield \eqref{eq:metric-transform}.  Finally,
\[
   g_i(y)(A_i(y)a,A_i(y)b)=B(a,b),
\]
so Proposition~\ref{prop:signature} transfers the signature.
\end{proof}

The important logical point is that the metric is \emph{derived after} soldering data are supplied.  It is not a field of the gluing datum and is not
inserted into the local input.

\subsection{Orientation compatibility}

Let
\[
   F_{ij}(y)
   :=
   A_j(\phi_{ij}(y))^{-1}
   \circ T_{ij}(y)\circ A_i(y).
\tag{8.3}
\]
By the intertwining identity \eqref{eq:soldering-law} this is the projected Lorentz transition (with the inverse-index convention induced by
\eqref{eq:soldering-law}).  Therefore
\[
   \det F_{ij}(y)=1.
\tag{8.4}
\]

\begin{proposition}[Orientation compatibility induced by soldering data]
\label{prop:orientation-reduction}
The local tangent frames obtained from the maps \(A_i\) have transition matrices
of determinant \(1\).  Equivalently, if \(a_i=\det A_i\), then on every overlap
\[
   \det T_{ij}(y)=\frac{a_j(\phi_{ij}(y))}{a_i(y)},
\]
so the multiplicative \Cech\ 1-cocycle formed by the determinants of the tangent transition maps is a continuous coboundary.
\end{proposition}

\subsection{Compatible future cones}

Define the interior future cone by
\[
   \mathcal C^{+}_{\mathrm{int}}
   :=
   \{x\in\Smodel : x\in C_{\Smodel}\ \text{and}\ N(x)>0\}.
\]
This uses only the already constructed cone of squares and quadratic
form.  In chart \(i\), define
\[
   \mathcal C^{+}_{i}(y)
   :=
   A_i(y)\mathcal C^{+}_{\mathrm{int}}.
\tag{8.5}
\]

\begin{proposition}[Future-cone compatibility]
\label{prop:future-cone}
For \(y\in W_{ij}\),
\[
   u\in\mathcal C^{+}_{i}(y)
   \quad\Longleftrightarrow\quad
   T_{ij}(y)u
   \in
   \mathcal C^{+}_{j}(\phi_{ij}(y)).
\tag{8.6}\label{eq:future-cone-transform}
\]
Thus the soldering data determine, in each tangent fibre, a continuous chart-independent choice of one of the two timelike cone components, namely the component corresponding to \(C^{+}_{\mathrm{int}}\).
\end{proposition}

The local vector field
\[
   \tau_i(y):=A_i(y)\mathbf1_{\Smodel}
\tag{8.7}
\]
is smooth, future-directed, and unit timelike.  It need not glue to a
single global field.  More precisely, the distinguished fields
\(\tau_i\) glue across an overlap exactly when the corresponding projected
Lorentz transition fixes \(\mathbf1_{\Smodel}\).

\begin{remark}[Scope: continuous choice of a timelike cone component]
\label{bdry:timeorientation}
The conclusion established here is the continuous compatible choice, in each tangent fibre, of the timelike cone component corresponding to \(C^{+}_{\mathrm{int}}\)
\eqref{eq:future-cone-transform}.  No theorem in this paper identifies this
choice with every standard equivalent formulation of time orientability.
Likewise, existence of an arbitrary global smooth future-directed timelike
section would require an additional global section theorem for the bundle
constructed here.
\end{remark}

% ======================================================================
\clearpage
\section{\Cech\ lifting obstruction for the central extension and tangent--Lorentz compatibility}
\label{sec:lifting-compatibility}
% ======================================================================

We now analyze the lifting problem for the Lorentz-valued cocycle on the chosen cover.

\subsection{Lifting obstruction on the chosen cover}

Let
\[
   g_{ij}:U_i\cap U_j\to\LorG
\]
be a Lorentz-valued \Cech\ 1-cocycle.  Assume continuous lifts
\[
   h_{ij}:U_i\cap U_j\to\SpinG,
   \qquad
   \rho(h_{ij})=g_{ij}.
\tag{9.1}
\]
have been chosen on the double overlaps.  The local-section theorem for
\(\rho\) guarantees such lifts only locally in the target; existence of a
continuous lift on an entire double overlap is additional information, or
may require a suitable refinement of the cover.  On triple overlaps, define
the defect
\[
   \epsilon_{ijk}
   :=
   h_{ij}h_{jk}h_{ik}^{-1}.
\tag{9.2}
\]
Because \(\rho(\epsilon_{ijk})=1\),
\[
   \epsilon_{ijk}\in\ker\rho=\{\pm1\}.
\tag{9.3}
\]
This defines a \(\{\pm1\}\)-valued \Cech\ 2-cocycle class on the chosen cover \(\cU\).
The independence from the chosen lift family is elementary at this level.
If \(h'_{ij}\) is another lift of the same Lorentz-valued cocycle, then
there are \(\{\pm1\}\)-valued functions
\(\alpha_{ij}:U_i\cap U_j\to\{\pm1\}\) with
\[
    h'_{ij}=\alpha_{ij}h_{ij}.
\]
Because the kernel is central,
\[
\begin{aligned}
 \epsilon'_{ijk}
 &=h'_{ij}h'_{jk}(h'_{ik})^{-1}\\
 &=\alpha_{ij}\alpha_{jk}\alpha_{ik}^{-1}\,
   \epsilon_{ijk}.
\end{aligned}
\tag{9.3a}\label{eq:obstruction-coboundary}
\]
Thus two representatives differ by the \(\{\pm1\}\)-valued \Cech\ coboundary \(\delta\alpha\), and determine the same \Cech\ 2-cocycle class on the chosen cover.
A \(\SpinG\)-valued \Cech\ 1-cocycle satisfying the cocycle law gives \(\epsilon_{ijk}=1\) identically, so
its class is trivial.

This is the familiar mechanism underlying the usual obstruction theory for
Spin structures \cite{Steenrod,MilnorStasheff,LawsonMichelsohn}; however,
the class considered here is defined only on the chosen cover \(\cU\).

\begin{remark}[Scope: no Stiefel--Whitney identification]
\label{bdry:no-sw}
No theorem here identifies the determinant condition derived from the transition maps on the chosen cover with \(w_1(TM)\),
or the \(\{\pm1\}\)-valued lifting class on the chosen cover with \(w_2(TM)\).
No refinement limit or comparison with global
\(H^{1}(M;\Ztwo)\) or \(H^{2}(M;\Ztwo)\) is used.  Consequently the
statement \(w_1(TM)=w_2(TM)=0\) is not a theorem of the present paper.
\end{remark}

\subsection{Intertwining under soldering data}

The intertwining identity \eqref{eq:soldering-law} identifies the projected Lorentz-valued cocycle with the tangent transition cocycle expressed in the local frames defined by \(A_i\).

\begin{theorem}[Intertwining of the tangent and Lorentz cocycles]
\label{thm:cocycle-identification}
For every choice of soldering data and every overlap, the change of tangent components between the local frames defined by \(A_i\) is exactly the Lorentz-valued transition \(g_{ij}\).
\end{theorem}

This is a rearrangement of the intertwining equation \eqref{eq:soldering-law}, but it is the precise intertwining relation between the two cocycles.

\begin{theorem}[Consequences of the intertwining identity]
\label{thm:soldering-consequences}
Assume \(B\) is smoothly glued and soldering data exist for a \(\SpinG\)-valued \Cech\ 1-cocycle \(\widetilde g\).  Then:
\begin{enumerate}
\item there are continuous nowhere-zero functions
      \(a_i:D_i\to\R\) such that
\[
  \det T_{ij}(y)\,a_i(y)
  =
  a_j(\phi_{ij}(y));
\tag{9.4}\label{eq:orientation-coboundary}
\]
\item the tangent transition cocycle in the local frames defined by \(A_i\) equals
      \(\rho(\widetilde g_{ij})\);
\item every local lift family of that projected cocycle has trivial \(\{\pm1\}\)-valued \Cech\ 2-cocycle class for the central extension on the chosen cover;
\item the compatible future timelike cone field of
      Proposition~\ref{prop:future-cone} exists.
\end{enumerate}
\end{theorem}

\begin{proof}
For (1), take the determinant of the intertwining identity \eqref{eq:soldering-law}.
The Lorentz-valued factor has determinant one, so the determinant of the
tangent transition is the ratio of the determinants of the local maps \(A_i\), giving
\eqref{eq:orientation-coboundary}.  Statement (2) is
\cref{thm:cocycle-identification}.  For (3), the \(\SpinG\)-valued \Cech\ 1-cocycle itself
is a coherent lift of its projection; the triple-overlap defect is
therefore trivial.  Independence of the obstruction from the chosen local
lift family transfers triviality to every lift family.  Statement (4) is
Proposition~\ref{prop:future-cone}.
\end{proof}

Two useful contrapositives are immediate.

\begin{corollary}[Necessary determinant condition for the existence of soldering data]
\label{cor:orientation-obstruction}
If no continuous nowhere-zero \(0\)-cochain \(a_i\) can satisfy
\eqref{eq:orientation-coboundary}, then no soldering data exist for any \(\SpinG\)-valued \Cech\ 1-cocycle over that gluing.
\end{corollary}

\begin{corollary}[Obstruction to lifting the Lorentz cocycle through the central extension]
\label{cor:spin-obstruction}
If the \(\{\pm1\}\)-valued \Cech\ 2-cocycle class associated with the central extension is nontrivial on the chosen cover,
then it admits no coherent \(\SpinG\)-valued \Cech\ 1-cocycle projecting to
it.  In particular, it cannot be the Lorentz-valued cocycle obtained from soldering data whose lift is a coherent \(\SpinG\)-valued \Cech\ 1-cocycle.
\end{corollary}

The converse of \cref{thm:soldering-consequences} is not proved.  Vanishing of these two necessary conditions is therefore not asserted to imply the existence of soldering data.

% ======================================================================
\clearpage
\section{Examples, obstructions, and nonuniqueness of lifts}
\label{sec:examples}
% ======================================================================

The following examples distinguish necessary conditions from features of
specific constructions.

\subsection{One-chart example}

On a one-chart quotient there are no nontrivial overlap constraints.  The gluing is smooth and soldering data exist.  Thus neither of the preceding obstructions on the chosen cover occurs in this example.

\subsection{Orientation-reversing counterexample}

There is a smooth orientation-reversing two-chart gluing whose cross-overlap has two
relevant components.  On the wrap-around component the tangent
transition has determinant \(-1\), whereas on the ordinary overlap component
it has determinant \(+1\).

\begin{proposition}[Orientation reversal prevents the existence of soldering data]
\label{prop:orientation-reversal}
For the orientation-reversing example, no soldering data exist for \emph{any} \(\SpinG\)-valued \Cech\ 1-cocycle.
\end{proposition}

\begin{proof}[Reason]
Every projected Lorentz-valued transition lies in \(\LorG\) and hence has
determinant \(+1\).  If soldering data existed, the determinant of each
local map \(A_i\) would be a continuous nowhere-zero function on a
connected chart domain and therefore would have constant sign there.  The determinant of the intertwining identity \eqref{eq:soldering-law} on the wrap-around component, where the tangent
determinant is \(-1\), forces opposite signs in the two charts; the same law
on the ordinary component, where the tangent determinant is \(+1\), forces
the same signs.  These requirements are incompatible.  This is the concrete
mechanism behind Corollary~\ref{cor:orientation-obstruction}; the
presence of both overlap components is essential.
\end{proof}

\(\SpinG\)-valued \Cech\ 1-cocycles nevertheless exist over this same gluing.
Thus existence of a \(\SpinG\)-valued \Cech\ 1-cocycle is strictly weaker than existence of soldering data intertwining its Lorentz projection with \(TM\).

\subsection{Combinatorial labels do not determine the existence of a soldering isomorphism}

Consider also a deliberately abstract combinatorial label structure.  For arbitrary such combinatorial labels,
the same label input is compatible with the one-chart example, where local
soldering data exist, and with the orientation-reversing two-chart example, where no soldering data exist for any \(\SpinG\)-valued \Cech\ 1-cocycle.  Therefore these combinatorial
labels do not determine the existence of soldering data.  

\subsection{Nonuniqueness of lifts modulo kernel-valued \Cech\ 0-cochains}

Because \(\ker\rho=\{\pm1\}\), distinct lifts of the same Lorentz-valued cocycle may remain inequivalent under \(\{\pm1\}\)-valued \Cech\ 0-cochains on the chosen cover.

\begin{proposition}[Two inequivalent lifts with identical Lorentz projection]
\label{prop:one-loop-freedom}
On the chosen one-loop cover there are two \(\SpinG\)-valued \Cech\ 1-cocycles that:
\begin{enumerate}
\item have the same projected Lorentz-valued cocycle;
\item both admit soldering data;
\item are not equivalent under the equivalence relation generated by \(\{\pm1\}\)-valued \Cech\ 0-cochains on the chosen cover.
\end{enumerate}
\end{proposition}

\begin{proposition}[Independent sign choices on a finite loop family]
\label{prop:multi-loop}
For \(N\) independent loop labels, the construction gives an injection of
\(\{\pm1\}^{N}\) into the equivalence classes under \(\{\pm1\}\)-valued \Cech\ 0-cochains on the chosen cover.  In particular, for
\(N=2\) there are four pairwise distinguished classes, all compatible with soldering data in the periodic gluing example.
\end{proposition}

\begin{proof}[Proof of Propositions~\ref{prop:one-loop-freedom} and~\ref{prop:multi-loop}]
For each loop \(n\), the cross-overlap of its two patches has two connected
components: the wrap-around component and the ordinary component.  Given a
sign assignment \(\sigma\in\{\pm1\}^{N}\), define a kernel-valued
\Cech~1-cocycle \(z^\sigma\) which equals \(\sigma_n\) on the
wrap-around component of loop \(n\) and equals \(1\) on the ordinary
component and on all unaffected overlaps.  Since every value of
\(z^\sigma\) lies in \(\ker\rho=\{\pm1\}\), twisting the trivial \(\SpinG\)-valued \Cech\ 1-cocycle by \(z^\sigma\) leaves the projected Lorentz cocycle unchanged.
Consequently the same soldering data for the periodic gluing example work for every
sign assignment.

It remains to prove that distinct sign assignments are not related by a
gauge transformation by \(\{\pm1\}\)-valued \Cech\ 0-cochains on the chosen cover.  Suppose \(\sigma_n\ne\tau_n\) for some
loop \(n\), and assume that patchwise kernel-valued gauge functions
\(\lambda_-\) and \(\lambda_+\) relate the corresponding twisted
cocycles.  The kernel is discrete and each patch is connected, so continuity
forces \(\lambda_-\) and \(\lambda_+\) to be constant on their patches.
Evaluating the gauge law at a point of the ordinary overlap, where both sign
cocycles equal \(1\), gives
\[
    \lambda_-\lambda_+^{-1}=1.
\]
Evaluating the same law at a point of the wrap-around component gives, using
centrality of the kernel,
\[
    \tau_n
      =\lambda_-\,\sigma_n\,\lambda_+^{-1}
      =\sigma_n,
\]
contradicting \(\sigma_n\ne\tau_n\).  Hence different sign assignments
give different equivalence classes under \(\{\pm1\}\)-valued \Cech\ 0-cochains on the chosen cover.  For one loop this yields exactly
the two classes of Proposition~\ref{prop:one-loop-freedom}; for \(N\)
independent loops it gives the asserted injection, and for \(N=2\) the four
sign assignments give four pairwise inequivalent classes.
\end{proof}

\begin{remark}[Scope]
The preceding statements do not compute \(H^{1}(M;\Ztwo)\).  They concern equivalence under \(\{\pm1\}\)-valued \Cech\ 0-cochains on the chosen cover; no refinement argument identifying these equivalence classes with global cohomology is part of the present construction.
\end{remark}

Hence the projected Lorentz-valued cocycle and soldering data do not determine a
unique lift to \(\SpinG\) on the chosen cover.

% ======================================================================
\clearpage
\section{A one-parameter family of \Cech\ 1-cocycles}
\label{sec:deformation}
% ======================================================================

We next study an explicit one-parameter family of \(\SpinG\)-valued \Cech\ 1-cocycles on a fixed periodic gluing.  No connection one-form, path-ordered
transport, holonomy functional, curvature two-form, or field equation is
introduced.

\subsection{A one-parameter subgroup}

Choose the first Euclidean basis direction \(\varepsilon_{0}\in V\)
and its Clifford image
\[
    \gamma_{0}:=\iota(\varepsilon_{0}),
    \qquad \gamma_{0}^{2}=1.
\]
The generator \(\gamma_{0}\) is chosen for this example; no theorem makes
this direction canonical.

\begin{definition}[One-parameter subgroup of \(\SpinG\)]
\label{def:G-lambda}
For \(\lambda\in\R\), define
\[
   G(\lambda)
   :=
   \cosh\!\left(\frac{\lambda}{2}\right)
   +
   \sinh\!\left(\frac{\lambda}{2}\right)\gamma_{0}
   \in\Cl_{3}.
\tag{11.1}\label{eq:selected-spin-path}
\]
\end{definition}

The hyperbolic identities and \(\gamma_{0}^{2}=1\) give
\[
   G(\lambda)c(G(\lambda))=1,
\]
hence \(G(\lambda)\in\SpinG\).  Moreover,
\[
   G(0)=1,\qquad
   G(\lambda+\mu)=G(\lambda)G(\mu),\qquad
   G(-\lambda)=G(\lambda)^{-1}.
\tag{11.2}
\]
The parameterization is injective.

The projection \(\rho(G(\lambda))\) is nontrivial for
\(\lambda\neq0\).  On the plane spanned by the scalar direction
and \((0,\varepsilon_{0})\), its action is the standard boost block
\[
   \begin{pmatrix}
      \cosh\lambda & \sinh\lambda\\
      \sinh\lambda & \cosh\lambda
   \end{pmatrix},
\tag{11.3}\label{eq:boost-block}
\]
and it acts trivially on the two orthogonal Euclidean directions.
In particular,
\[
   \rho(G(\lambda))\mathbf1_{\Smodel}
   =
   \bigl(\cosh\lambda,\sinh\lambda\,\varepsilon_{0}\bigr).
\tag{11.4}
\]

\begin{remark}[Scope of the example]
\label{bdry:selected}
The family \eqref{eq:selected-spin-path} is a deliberately selected one-parameter subgroup
generated by one chosen Clifford direction.  It is neither proved unique
nor canonical, and it does not exhaust other one-parameter subgroups or cocycle families.
\end{remark}

\subsection{The periodic gluing example}

Fix the periodic gluing, together with the chosen cover and two-component periodic overlap.  We make no identification of
this quotient with \(S^{1}\times\R^{3}\), and no fundamental-group theorem
is required.

The \(\SpinG\)-valued \Cech\ 1-cocycle is varied using \(G(\lambda)\).
The gluing and hence the tangent transitions are held fixed; the
projected Lorentz-valued transition is obtained by applying \(\rho\) to the \(\SpinG\)-valued transition.  We then ask for local soldering data satisfying the intertwining identity \eqref{eq:soldering-law}.

Let
\[
  \Lambda_{\mathrm{sold}}
  :=
  \left\{
      \lambda\in\R:
      \text{soldering data exist for the cocycle at }\lambda
  \right\}.
\tag{11.5}
\]

\begin{theorem}[Existence of soldering data for every parameter]
\label{thm:admissibility}
For the fixed periodic gluing,
\[
   \Lambda_{\mathrm{sold}}=\R.
\tag{11.6}\label{eq:admissibility-region}
\]
Thus every finite parameter value admits explicitly constructed soldering data.  In particular existence is symmetric under
\(\lambda\mapsto-\lambda\).
\end{theorem}

\begin{proof}
Let \(\chi:\R\to[0,1]\) be a fixed \(C^\infty\) transition profile with
\(\chi(t)=0\) for \(t\le3\) and \(\chi(t)=1\) for \(t\ge4\).  On the two
patches of each periodic loop define
\[
  \alpha_\lambda(p,y)=
  \begin{cases}
     0, & p\text{ is the first patch},\\
     \lambda\,\chi(y_0), & p\text{ is the second patch}.
  \end{cases}
\tag{11.6a}\label{eq:solder-angle}
\]
Let \(P(t)=\rho(G(t))\in\LorG\).  By the one-parameter subgroup law,
\(P(s+t)=P(s)P(t)\) and \(P(-t)=P(t)^{-1}\), and the closed form
\eqref{eq:boost-block} makes \(t\mapsto P(t)\) smooth.  Define
\[
   A_p^{(\lambda)}(y):=P(\alpha_\lambda(p,y)).
\tag{11.6b}\label{eq:explicit-solder}
\]
This is the identity on the first patch and smoothly interpolates on the
second patch between the identity and \(P(\lambda)\).

All tangent transitions of the fixed periodic gluing are the identity.
The projected Lorentz-valued transition on an overlap is
\(P(\beta_{pq}^{(\lambda)}(x))\), where
\(\beta_{pq}^{(\lambda)}\) is \(\pm\lambda\) on the wrap-around component
and \(0\) on the ordinary component.  By the defining values of \(\chi\),
one verifies on each overlap component that
\[
   \alpha_\lambda(q,\phi_{pq}(y))
   =\alpha_\lambda(p,y)+\beta_{pq}^{(\lambda)}(x_y).
\tag{11.6c}\label{eq:solder-angle-step}
\]
Therefore
\[
\begin{aligned}
  A_q^{(\lambda)}(\phi_{pq}(y))
   &=P(\alpha_\lambda(q,\phi_{pq}(y)))\\
   &=P(\alpha_\lambda(p,y))
     P(\beta_{pq}^{(\lambda)}(x_y))\\
   &=A_p^{(\lambda)}(y)\,g_{pq}^{(\lambda)}(x_y),
\end{aligned}
\]
which is precisely the intertwining identity \eqref{eq:soldering-law} because the tangent transition is
\(\id\).  Smoothness and pointwise invertibility follow from those of
\(P\) and \(\chi\).  Hence soldering data exist for every
\(\lambda\in\R\), proving \(\Lambda_{\mathrm{sold}}=\R\).
\end{proof}

\subsection{Gauge equivalence by group-valued \Cech\ 0-cochains}

Let \(\widetilde g^{(\lambda)}\) denote the \(\SpinG\)-valued \Cech\ 1-cocycle at
parameter \(\lambda\).

\begin{theorem}[Gauge equivalence of the one-parameter family]
\label{thm:gauge-orbit}
For every pair \(\lambda_{1},\lambda_{2}\in\R\) there are continuous maps \(u_i:U_i\to\SpinG\) such that, on each overlap,
\[
   \widetilde g_{ij}^{(\lambda_2)}
   =u_i\,\widetilde g_{ij}^{(\lambda_1)}\,u_j^{-1}.
\tag{11.7}
\]
The \(\SpinG\)-valued \Cech\ 0-cochain \(u\) is continuous.  After projection by \(\rho\), the resulting
Lorentz-valued \Cech\ 0-cochain is \(C^\infty\) on every chart domain, and the
projected Lorentz cocycles are gauge equivalent by that projected 0-cochain.
\end{theorem}

\begin{proof}
First compare an arbitrary \(\lambda\) with the identity member \(\lambda=0\).  Using the
same interpolation function \(\alpha_\lambda\) as in
\eqref{eq:solder-angle}, define on patch \(p\)
\[
   u_p^{(\lambda)}(x)
   :=G\!\left(-\alpha_\lambda
       \bigl(p,\operatorname{coord}_p(x)\bigr)\right).
\tag{11.7a}\label{eq:spin-gauge}
\]
The map is continuous because \(G(t)\) is continuous in the topology on \(\SpinG\) and the chart-coordinate profile is continuous.  On an overlap,
\eqref{eq:solder-angle-step} and the group law of \(G\) give
\[
   \widetilde g_{pq}^{(\lambda)}(x)
   =u_p^{(\lambda)}(x)\,
      \widetilde g_{pq}^{(0)}(x)\,
      \bigl(u_q^{(\lambda)}(x)\bigr)^{-1}.
\tag{11.7b}\label{eq:spin-coboundary}
\]
Thus every parameter value is gauge equivalent under \(\SpinG\)-valued \Cech\ 0-cochains to the identity cocycle, i.e. the member with \(\lambda=0\).
For two arbitrary values \(\lambda_1,\lambda_2\), compose the inverse gauge
from \(\lambda_1\) to the identity member with the gauge from the identity member to \(\lambda_2\).

Applying \(\rho\) to \eqref{eq:spin-coboundary} gives the corresponding
Lorentz gauge relation.  Its chartwise representative is
\(P(-\alpha_\lambda)\); by the explicit hyperbolic formula
\eqref{eq:boost-block} and smoothness of \(\alpha_\lambda\), this projected
gauge is \(C^\infty\) on every chart domain.  No smooth Lie-group structure on \(\SpinG\) is used or claimed.
\end{proof}

The equivalence relation generated only by \(\{\pm1\}\)-valued \Cech\ 0-cochains on the chosen cover behaves differently.
For \(\lambda\neq0\), the cocycle at parameter \(\lambda\) is not obtained from the identity cocycle merely by a \(\{\pm1\}\)-valued
0-cochain, because its Lorentz projection is nontrivial.  There is no contradiction: the two equivalence relations answer
different questions.

\subsection{Bijection between sets of local soldering data}

For a parameter \(\lambda\), write
\[
   \cE_{\lambda}
   :=
   \left\{
      \text{soldering data for }
      (B_{\mathrm{per}},\widetilde g^{(\lambda)})
   \right\}.
\tag{11.8}
\]

\begin{theorem}[Bijection between sets of local soldering data]
\label{thm:solution-correspondence}
For every \(\lambda_{1},\lambda_{2}\in\R\), there is an explicit bijection
\[
   \mathcal T_{\lambda_{1},\lambda_{2}}:
   \cE_{\lambda_{1}}
   \xrightarrow{\;\sim\;}
   \cE_{\lambda_{2}},
\tag{11.9}\label{eq:solution-bijection}
\]
obtained by precomposing each local representative \(A_i\) with the projected
gauge relating the two \(\SpinG\)-valued \Cech\ 1-cocycles.  The inverse is
\(\mathcal T_{\lambda_{2},\lambda_{1}}\).
\end{theorem}

\begin{proof}
Set
\[
   \delta_{12}(p,y)
   :=\alpha_{\lambda_2}(p,y)-\alpha_{\lambda_1}(p,y),
   \qquad
   Q_{12,p}(y):=P(\delta_{12}(p,y)).
\tag{11.9a}\label{eq:gauge-angle}
\]
For soldering data \(A\in\cE_{\lambda_1}\), define
\[
   \bigl(\mathcal T_{\lambda_1,\lambda_2}A\bigr)_p(y)
   :=A_p(y)\circ Q_{12,p}(y).
\tag{11.9b}\label{eq:solder-transport}
\]
The fields \(Q_{12,p}\) are smooth and invertible.  Subtracting the two
identities \eqref{eq:solder-angle-step} for \(\lambda_1\) and
\(\lambda_2\) gives exactly the overlap equation needed to convert the
the intertwining identity at \(\lambda_1\) into the intertwining identity at \(\lambda_2\).  Thus
\eqref{eq:solder-transport} is soldering data at \(\lambda_2\).

Since
\[
   \delta_{21}=-\delta_{12}
   \quad\text{and}\quad
   P(-t)=P(t)^{-1},
\]
transporting from \(\lambda_1\) to \(\lambda_2\) and then back composes
with \(Q_{12}Q_{21}=\id\) pointwise.  Hence
\(\mathcal T_{\lambda_2,\lambda_1}\) is the inverse of
\(\mathcal T_{\lambda_1,\lambda_2}\), proving the stated bijection.
\end{proof}

\begin{theorem}[Metric preservation]
\label{thm:metric-preservation}
Let \(A\in\cE_{\lambda_{1}}\), and let
\(A'=\mathcal T_{\lambda_{1},\lambda_{2}}(A)\).
Then for every \(x\in M\) and \(X,Y\in T_xM\),
\[
   g_{A'}(X,Y)=g_A(X,Y).
\tag{11.10}
\]
\end{theorem}

\begin{proof}
In a chart, \eqref{eq:solder-transport} gives
\[
    A'_p=A_p\circ Q_{12,p},
    \qquad
    (A'_p)^{-1}=Q_{12,p}^{-1}\circ A_p^{-1}.
\]
Therefore
\[
\begin{aligned}
 g_{A',p}(u,v)
 &=B\!\left(Q_{12,p}^{-1}A_p^{-1}u,
            Q_{12,p}^{-1}A_p^{-1}v\right)\\
 &=B\!\left(A_p^{-1}u,A_p^{-1}v\right)\\
 &=g_{A,p}(u,v),
\end{aligned}
\]
because \(Q_{12,p}(y)=P(\delta_{12}(p,y))\in\LorG\) preserves \(B\).
The equality holds in every chart and is compatible with the tangent
transition law, hence it is the equality of the induced Lorentzian metrics at
every base point and on every pair of tangent vectors.
\end{proof}

\begin{corollary}[Gauge equivalence and equality of the induced Lorentzian metrics]
\label{cor:family-gauge-redundancy}
All members of the one-parameter family are gauge equivalent under \(\SpinG\)-valued \Cech\ 0-cochains.  Under the explicit bijections between the sets of local soldering data,
corresponding choices of soldering data induce the same Lorentzian metric.
\end{corollary}

\begin{remark}[Scope of Corollary~\ref{cor:family-gauge-redundancy}]
The result does not classify arbitrary deformations of \(\SpinG\)-valued \Cech\ 1-cocycles, arbitrary
covers, co-varying bases, or connection-level deformations.  It does not
imply gauge equivalence for arbitrary cocycle families, and it carries no
curvature interpretation.  The parameter \(\lambda\) is the coordinate of
the selected one-parameter subgroup of \(\SpinG\) \eqref{eq:selected-spin-path}, not a curvature scalar.
\end{remark}

% ======================================================================
\clearpage
\section{Scope and open problems}
\label{sec:boundary}
% ======================================================================

The results above determine the following scope and open problems.

From the local input \(V=\R^{3}\) with its Euclidean form, one has
constructed the real spin factor with Lorentzian quadratic form, the Euclidean Clifford algebra, the two-fold central covering \(\SpinG\to\LorG\), and the six-dimensional orthogonal Lie algebra \(\mathfrak{so}(B)\).
After adding explicit gluing data and sufficient topological and smooth
conditions, one obtains a smooth four-manifold.  After adding soldering data, one obtains a smooth Lorentzian metric, orientation compatibility, a continuous compatible choice of the timelike cone component corresponding to \(C^{+}_{\mathrm{int}}\) in each tangent fibre, and the exact intertwining of the tangent transition cocycle with the projected Lorentz-valued cocycle.

The following structures and identifications have not been constructed in this paper.

\begin{enumerate}
\item The topological group \(\SpinG\) has not been given a smooth Lie-group
      structure and related to \(\soB\) as a Lie group/Lie algebra pair:
\[
   \operatorname{Lie}(\SpinG)\simeq\soB
   \qquad\text{is open.}
\tag{12.1}
\]
\item The differential \(d\rho_e\) of the covering homomorphism has not been
      constructed.
\item No connection data associated with either the \(\SpinG\)-valued or the Lorentz-valued \Cech\ cocycle have been introduced.
\item Consequently there is no covariant derivative, parallel transport,
      holonomy, curvature, torsion, action functional, or field equation in
      the present construction.
\item The determinant \Cech\ 1-cocycle and the \(\{\pm1\}\)-valued \Cech\ 2-cocycle class on the chosen cover have not been identified
      with \(w_1(TM)\) and \(w_2(TM)\), respectively.
\item The compatible future timelike cone field constructed in Proposition~\ref{prop:future-cone}
      has not been identified by theorem with all standard formulations of
      time orientability.
\end{enumerate}

In the periodic gluing example of \cref{sec:deformation}, the entire one-parameter
family is gauge equivalent under \(\SpinG\)-valued \Cech\ 0-cochains, and the induced Lorentzian metric is unchanged under the explicit bijection between the sets of local soldering data.  Thus this particular transition-cocycle parameter does not define
distinct induced Lorentzian metrics at the level studied here.  No conclusion about
connection-level geometry follows from this example.

% ======================================================================
\clearpage
\section{Conclusion}
% ======================================================================

The results establish a separation between levels of geometric
information.

At the local level, from the specified positive Euclidean three-space one constructs the real spin factor \(\Smodel\) with its Lorentzian quadratic form, the Euclidean Clifford algebra, the two-fold central covering \(\SpinG\to\LorG\), and the orthogonal Lie algebra \(\mathfrak{so}(B)\).  This part is determined by the local input.

At the global level, gluing information is independent.  Local pieces,
even when equipped with fibre transitions, do not determine which points
of different pieces represent the same point of the quotient.  Explicit gluing data are therefore additional input.  When the graph of the gluing equivalence relation is closed, the index set is countable, and the transition maps are smooth, the quotient is a smooth four-manifold.

At the tangent-bundle level, a second independent datum is required.  The derivative cocycle \(D\phi_{ij}\) of that manifold is not forced to equal the independently specified Lorentz-valued cocycle.  Local soldering data provide the vector-bundle isomorphism \(s:E_g\to TM\) intertwining the associated Lorentz vector bundle with \(TM\).  These data yield a smooth Lorentzian metric, orientation compatibility, a continuous compatible choice of the timelike cone component corresponding to \(C^{+}_{\mathrm{int}}\) in each tangent fibre, and triviality, on the chosen cover, of the \(\{\pm1\}\)-valued \Cech\ 2-cocycle class associated with the central extension.

The resulting implication is therefore conditional and exact:
\[
\begin{array}{c}
(\R^{3},\langle\cdot,\cdot\rangle)
\Longrightarrow
(\Smodel,N,\Cl_{3},\SpinG\!\to\!\LorG,\soB)
\\[0.5em]
\begin{gathered}
\text{local algebraic data}+\text{gluing}\\
+\text{closed gluing relation}+\text{countable index}+\text{smoothness}
\end{gathered}
\Longrightarrow M^{4}
\\[0.5em]
M^{4}+\text{\(G_{\mathrm{Spin}}\)-valued \Cech\ 1-cocycle}+\text{soldering data}
\\[-0.1em]
\Downarrow
\\[-0.1em]
(TM,g),
\\
\text{orientation compatibility},
\\
\text{compatible choice of a timelike cone component in each tangent fibre},
\\
\text{tangent cocycle in the local frames defined by }A_i\text{ = projected Lorentz cocycle},
\\
\text{trivial \(\{\pm1\}\)-valued \Cech\ 2-cocycle class for the central extension on the chosen cover}.
\end{array}
\tag{13.1}\label{eq:final-chain}
\]
Each implication in \eqref{eq:final-chain} is conditional on the displayed
hypotheses.

In the periodic gluing example of \cref{sec:deformation}, the selected
members of the one-parameter family are gauge equivalent for all finite parameters.  The
sets of local soldering data at different parameter values are explicitly
bijective, and corresponding choices of soldering data induce the same Lorentzian metric.  This
result is restricted to the specified family and does not address
connection-level geometry.

% ======================================================================
\appendix
\clearpage
\section{Dependency and status table}
\label{app:status}
% ======================================================================

\begin{center}
\small
\begin{tabular}{>{\raggedright\arraybackslash}p{0.30\textwidth}
                >{\raggedright\arraybackslash}p{0.23\textwidth}
                >{\raggedright\arraybackslash}p{0.37\textwidth}}
\toprule
Object or statement & Status & Dependency \\
\midrule
\(V=\R^{3}\), Euclidean form & input & none \\
\(\Smodel,N,B\), signature \((1,3)\) & \statusderived & Euclidean input \\
\(\Cl_{3}\), paravectors & \statusderived & \(q_{3}\) \\
\(\SpinG\to\LorG\), kernel \(\{\pm1\}\) & \statusderived & Clifford algebra and paravector action \\
\(\soB\simeq\Lambda^{2}\Smodel\) & \statusderived & \(B\) \\
Gluing data & \statusdatum & not determined by local pieces \\
Closed graph of the gluing equivalence relation & additional hypothesis & used in the Hausdorff proof \\
Countable index set & additional hypothesis & used in the second-countability proof \\
Smooth gluing & sufficient additional hypothesis & base transitions \(\phi_{ij}\) \\
Smooth \(M^{4}\) & \statusderived & preceding global hypotheses \\
Tangent transitions \(D\phi_{ij}\) & \statusderived & smooth atlas \\
\(\SpinG\)-valued \Cech\ 1-cocycle & global datum & chosen cover of \(M\) \\
Soldering data & \statusdatum & existence not automatic \\
Smooth Lorentzian metric & \statusderived & soldering data \\
Orientation compatibility & \statusderived & soldering data \\
Compatible choice of a timelike cone component in each tangent fibre & \statusderived & soldering data \\
\(\{\pm1\}\)-valued \Cech\ 2-cocycle class trivial on the chosen cover & \statusderived & soldering data + \(\SpinG\)-valued \Cech\ 1-cocycle \\
\(w_1,w_2\) comparison & \statusopen & characteristic-class comparison theorem \\
\(\operatorname{Lie}(\SpinG)\simeq\soB\) & \statusopen & smooth Lie-group structure on \(\SpinG\) \\
\(d\rho_e\) & \statusopen & preceding infinitesimal bridge \\
Connection, holonomy, curvature & \statusopen & additional connection data \\
\bottomrule
\end{tabular}
\end{center}

% ======================================================================
\clearpage
\section{Source correspondence}
\label{app:source-map}
% ======================================================================

This appendix gives the correspondence between the mathematical statements
in the paper and the Lean source.  Implementation identifiers retain the
historical project naming; these names are not mathematical terminology used
in the paper.

\begin{center}
\footnotesize
\begin{tabular}{>{\raggedright\arraybackslash}p{0.28\textwidth}
                >{\raggedright\arraybackslash}p{0.64\textwidth}}
\toprule
Paper statement & Principal Lean declaration(s) \\
\midrule
Spin factor and Lorentz form &
\path{SpinCore.sJ_jordan},
\path{SpinCore.lorentz_quadratic_form} \\
Jordan--Clifford bridge &
\path{SpinCore.clifford_jordan},
\path{SpinCore.adjoin_spinToCl_eq_top} \\
Two-fold central covering &
\path{SpinCore.spin_double_cover_bundled},
\path{SpinCore.spin_double_cover_topological_endpoint},
\path{SpinCore.spin_double_cover_local_section_endpoint} \\
Bivectors and \(\mathfrak{so}(B)\) &
\path{SpinCore.bivectorEquivSkew},
\path{SpinCore.finrank_gB},
\path{SpinCore.gBLie_bracket} \\
Local pieces do not determine the quotient space &
\path{EmergentBase.base_not_determined_by_local_pieces} \\
Topological quotient construction &
\path{EmergentBase.BaseGluingData.emergentManifoldCertificate} \\
Smooth manifold &
\path{EmergentBase.BaseGluingData.isManifold_of_smoothGluing},
\path{EmergentBase.BaseGluingData.smoothEmergentManifoldCertificate} \\
Tangent transition &
\path{EmergentBase.BaseGluingData.tangentTransition_eq_derivative_baseTransition} \\
Independence of tangent and Lorentz-valued cocycles &
\path{SpinNative.tangent_spin_base_data_do_not_force_transition_identification} \\
Soldering data and vector-bundle isomorphism &
\path{SpinNative.SmoothTangentSolderData};
\path{SpinNative.smooth_solder_iff_regular_bundle_equivalence} \\
Smooth Lorentzian metric &
\path{SpinNative.SmoothTangentSolderData.smooth_tangent_lorentz_metric} \\
Orientation compatibility and compatible timelike cone choice &
\path{SpinNative.SmoothTangentSolderData.orientation_time_orientation_status} \\
Intertwining of tangent and Lorentz-valued cocycles &
\path{Task36.classical_reconvergence},
\path{Task36.spin_obstruction_nonzero_implies_no_solder} \\
Counterexamples to existence of soldering data &
\path{Task36.adversarial_global_topology_certificate},
\path{Task36.spin_foam_labels_do_not_determine_solder} \\
One-parameter \(\SpinG\)-valued \Cech\ 1-cocycle family &
\path{Task37.Deformation.transportState_add},
\path{Task37.Deformation.projectedTransportState_apply} \\
Existence of soldering data for the periodic family &
\path{Task37.Deformation.loopParaSolder} (construction),
\path{Task37.Deformation.controlAAdmissible_all},
\path{Task37.Deformation.regularRegion_loopTransportFamily} \\
Gauge equivalence under \(\SpinG\)-valued \Cech\ 0-cochains &
\path{Task37.Deformation.loopGaugeFun} (construction),
\path{Task37.Deformation.spinCocycle_gaugeEquiv},
\path{Task37.Deformation.projectedCocycle_gaugeEquiv} \\
Combined theorem for the one-parameter family &
\path{Task37.Deformation.native_transport_knob_smoketest_certificate} \\
Bijection between sets of local soldering data &
\path{Task37.Deformation.loopSolderEquiv} (construction),
\path{Task37.Deformation.loopSolderTransport_involutive},
\path{Task37.Deformation.loopSolderSolutionSpace_correspondence},
\path{Task37.Deformation.tangentMetric_loopSolderTransport} \\
\bottomrule
\end{tabular}
\end{center}

The formal development is available at
\url{https://github.com/antaris82/The_Relational_Spin-Closure_Hypothesis}
\cite{RSCHFormal}.
The main development is written in Lean~4 and uses Mathlib
\cite{LeanSystem,Mathlib2020}.  The mathematical claims of this paper should
be read at the strength of the displayed hypotheses above; the Lean source supplies machine-checked proofs of the corresponding
statements.

% ======================================================================
\clearpage
\section*{Acknowledgments}
\addcontentsline{toc}{section}{Acknowledgments}
% ======================================================================

The editor thanks Thomas Stoer and Thomas Klimpel for sustained, critical,
and constructive discussions.  These exchanges repeatedly helped clarify
assumptions, expose weak formulations, and sharpen the mathematical
questions addressed in this work.

% ======================================================================
\section*{AI and formalization disclosure}
\addcontentsline{toc}{section}{AI and formalization disclosure}
% ======================================================================

ChatGPT (OpenAI), acting as AI Assistant, was used for mathematical
exploration, source-level analysis, organization, consistency checking, and
editorial drafting.  Aristotle was used as Lean Agent for prover-assisted
formalization and verification.  The editor curated the scope and the
presentation.  Mathematical claims are asserted only at the strength stated
in the paper and are tied, where applicable, to the corresponding Lean declarations.  The Lean source contains the machine-checked proofs of those statements.

% ======================================================================
\small
\setlength{\emergencystretch}{2em}
\begin{thebibliography}{99}

\bibitem{Chevalley}
C.~Chevalley,
\emph{The Algebraic Theory of Spinors},
Columbia University Press, New York, 1954.

\bibitem{FarautKoranyi}
J.~Faraut and A.~Kor\'anyi,
\emph{Analysis on Symmetric Cones},
Oxford Mathematical Monographs,
Oxford University Press, Oxford, 1994.
doi:10.1093/oso/9780198534778.001.0001.

\bibitem{Husemoller}
D.~Husemoller,
\emph{Fibre Bundles},
3rd ed., Graduate Texts in Mathematics 20,
Springer, New York, 1994.

\bibitem{JordanNeumannWigner}
P.~Jordan, J.~von Neumann, and E.~Wigner,
On an algebraic generalization of the quantum mechanical formalism,
\emph{Annals of Mathematics} \textbf{35} (1934), 29--64.
doi:10.2307/1968117.

\bibitem{KobayashiNomizu}
S.~Kobayashi and K.~Nomizu,
\emph{Foundations of Differential Geometry, Vol.~I},
Interscience Publishers, New York, 1963.

\bibitem{LawsonMichelsohn}
H.~B.~Lawson, Jr. and M.-L.~Michelsohn,
\emph{Spin Geometry},
Princeton Mathematical Series 38,
Princeton University Press, Princeton, 1989.

\bibitem{LeeSmooth}
J.~M.~Lee,
\emph{Introduction to Smooth Manifolds},
2nd ed., Graduate Texts in Mathematics 218,
Springer, New York, 2013.

\bibitem{MilnorStasheff}
J.~W.~Milnor and J.~D.~Stasheff,
\emph{Characteristic Classes},
Annals of Mathematics Studies 76,
Princeton University Press, Princeton, 1974.

\bibitem{Nakahara}
M.~Nakahara,
\emph{Geometry, Topology and Physics},
2nd ed., CRC Press, Boca Raton, 2003.
ISBN 978-0-7503-0606-5.

\bibitem{ONeill}
B.~O'Neill,
\emph{Semi-Riemannian Geometry with Applications to Relativity},
Pure and Applied Mathematics 103,
Academic Press, New York, 1983.

\bibitem{Steenrod}
N.~Steenrod,
\emph{The Topology of Fibre Bundles},
Princeton Mathematical Series 14,
Princeton University Press, Princeton, 1951.

\bibitem{LeanSystem}
L.~de Moura, S.~Kong, J.~Avigad, F.~van Doorn, and J.~von Raumer,
The Lean theorem prover (system description),
in A.~P.~Felty and A.~Middeldorp (eds.),
\emph{Automated Deduction -- CADE-25},
Lecture Notes in Computer Science 9195,
Springer, Cham, 2015, pp.~378--388.

\bibitem{Mathlib2020}
The mathlib Community,
The Lean mathematical library,
in \emph{Proceedings of the 9th ACM SIGPLAN International Conference
on Certified Programs and Proofs (CPP 2020)},
ACM, 2020, pp.~367--381.
doi:10.1145/3372885.3373824.

\bibitem{RSCHFormal}
{\raggedright
\emph{Formal Lean development accompanying the construction results},
source repository, commit
\path{e89832ec0e1f301c051b01a0a4118d9082118a60},
\url{https://github.com/antaris82/The_Relational_Spin-Closure_Hypothesis},
accessed 12 September 2026.\par}

\end{thebibliography}

\end{document}
